\documentclass[11pt]{amsart}
\usepackage[margin=1in]{geometry}
\usepackage{amsmath,amssymb,amsthm,mathtools,booktabs}
\usepackage{enumitem}
\usepackage{hyperref}
\hypersetup{colorlinks=true,linkcolor=blue,citecolor=blue,urlcolor=blue}
\usepackage[final]{microtype}
\setlength{\emergencystretch}{3em}
\usepackage{url}
\urlstyle{same}

\newcommand{\leanrepo}{https://github.com/HQIV/hqiv-lean/blob/main}
\newcommand{\leanfile}[1]{\href{\leanrepo/#1}{\path{#1}}}
\newcommand{\leanid}[1]{\path{#1}}

\title[Regional $\zeta(s)$ and SO(4) projection]
{Regional Closed Forms for $\zeta(s)$:\\
An SO(4) Projection and Machine-Checked Reformulations of RH}

\author{Steven Ettinger Jr}
\address{Independent Researcher}
\email{steven@disregardfiat.tech}
\date{Draft --- Lean formal note (2026-06)}

\newtheorem{definition}{Definition}[section]
\newtheorem{proposition}[definition]{Proposition}
\newtheorem{theorem}[definition]{Theorem}
\newtheorem{lemma}[definition]{Lemma}
\newtheorem{remark}[definition]{Remark}
\newtheorem{corollary}[definition]{Corollary}

\begin{document}
\begin{abstract}
\noindent\textbf{Position.}
This note records a machine-checked geometric interpretation of the Riemann zeta
function in the public HQIV Lean library \cite{hqiv-lean}, building on the
$\mathfrak{so}(8)$ closure record \cite{hqiv-so8-paper}, the octonion
light-cone carrier \cite{hqiv-lightcone-paper}, and the TUFT+SM synthesis
\cite{hqiv-tuft-sm-lagrangian-paper}.  It is a \emph{formal mathematics} paper,
not a phenomenology forecast: every theorem is a regional closed form,
projection identity, or logical equivalence stated in standard analytic
notation and verified by \texttt{lake build HQIVStory}.

\noindent\textbf{Main geometric claim.}
The functional-equation reflection $\sigma\leftrightarrow 1-\sigma$ is realized
as a $45^\circ$ rotation in a quaternionic plane.  The fixed diagonal component
is $1/\sqrt2$; the variable \emph{equator factor} $(2\sigma-1)/\sqrt2$
(Lean: \texttt{so4CriticalFactor}) is the sin--cos carrier on the strip.
Classical $\zeta(s)$ is then a \emph{piecewise} closed form on $\mathbb{C}$ in
the usual sense: Dirichlet/Euler for $\Re s>1$, the functional-equation
sin--cos--$\Gamma$ assembly on the open strip, Bernoulli values at negative
integers, and $\pi$-power formulas at even positives.

\noindent\textbf{Open-strip assembly.}
The functional-equation expression for $\zeta(s)$ on $0<\Re s<1$ combines
$\Gamma(1-s)$ poles, oscillating sin/cos factors, and apparent singularities
that obscure where the \emph{nontrivial} zeros can live until the regional
bookkeeping is separated (\S\ref{sec:real-line}): trivial zeros and $\Gamma$
poles are classified explicitly; the equator factor vanishes only on
$\Re s=\tfrac12$; and the meaningful zero question is pinned to that critical
line.  Pairwise cancellation along the FE orbit $(s,1-s)$ holds for every
$\sigma$, but only pointwise vanishing of the equator factor forces
$\sigma=\tfrac12$ (\S\ref{sec:orbit-gap}).

\noindent\textbf{Sine/cosine trichotomy.}
At even integers $s=2k$ the angle $\pi s/2=k\pi$ kills the sine factor
($\sin=0$, $\cos=(-1)^k$): even $\zeta$-values are pure $\pi$-powers.
At odd integers the cosine factor vanishes and sine carries the residue (no
elementary closed form for $\zeta(3)$, $\zeta(5)$, \ldots).
On the fractional open strip both factors are nonzero and
$\tan(\pi s/2)$ is the rotation ratio.

\noindent\textbf{RH and Goldbach.}
Classically RH and Goldbach are independent open problems.  Here both are tied
to the same normalized slope $\tfrac12$: critical-line localization for
nontrivial $\zeta$-zeros, and midpoint slope $N/(p+q)=\tfrac12$ for Goldbach
prime pairs (\S\ref{sec:rh-goldbach}).  Lean proves an \emph{exact equivalence}
between a single bridge hypothesis (\texttt{SO8ProjectedHalfSlopeBridge}\,$2$)
and \texttt{RiemannHypothesis} $\wedge$ \texttt{GoldbachParity}---not a
one-sided implication with extra conditions.  Proving the bridge remains as hard
as the conjunction of the two classical conjectures.

\noindent\textbf{Status.}
Every algebraic / Mathlib link in the harmonic--prime--$\zeta$ chain, every
45$^\circ$ projection identity, the $\Re s$-dependence lemmas, the FE
quadruplet orbit formulation, the bridge equivalence
RH $\wedge$ Goldbach $\Leftrightarrow$ half-slope bridge, the
discrete-dominates-continuum heat-flow comparison with exact phase lock,
the concrete-$\Lambda$ port of the de Bruijn--Newman constant, the tangent
orbit law, the spectral-cancellation profile of zeros (no cancellation in
the Euler product region; resonance blow-up of $1/\zeta$ at every zero), and
the resonance-channeling reduction (two-channel dichotomy;
\texttt{InteriorAssemblyNonzero} $\Leftrightarrow$ RH; RH $\Leftrightarrow$
square-root spectral weights at every zero), the vertical-axis rigidity
(the zero set admits at most one mirror axis, necessarily $\Re s=\tfrac12$),
the Diophantine transformer laws (trace ladder $\to\zeta$; FE
composition $=\mathrm{diag}(1/n)$; adjoint law $\Leftrightarrow$ line; RH
$\Leftrightarrow$ adjoint law at every zero), and the octonionic
associator channel (commutative saturation of the transformer family;
non-associative torsion provably active at every zero; RH
$\Leftrightarrow$ torsion weights), and the S$^7$/SO(8) torsion
cancellation laws (single $S^7$ point; chirality flip; RH
$\Leftrightarrow$ FE torsion cancellation at every zero), and the log
edge (polar decoupling; prime-log incommensurability; two prime phases
pin the height; the Goldbach circle $pq+(q-N)^2=N^2$)
cited here are sorry-free in \texttt{lake build HQIVStory}.  Unconditional
\emph{Riemann Hypothesis} and \emph{Goldbach parity} are \emph{not} claimed: the
main conditional hypothesis
\texttt{InteriorAssemblyNonzeroAtNontrivialZerosOffLine} and a witness for the
(provably equivalent) half-slope bridge remain open (Appendix~\ref{app:status}).
\end{abstract}
\maketitle

% Scope note for Lean formal / mathematical notes (not full patch phenomenology).
% From papers/*.tex: % Scope note for Lean formal / mathematical notes (not full patch phenomenology).
% From papers/*.tex: % Scope note for Lean formal / mathematical notes (not full patch phenomenology).
% From papers/*.tex: \input{../include/readout_dictionary_messaging.tex}

\providecommand{\leanfile}[1]{\nolinkurl{#1}}
\providecommand{\leanid}[1]{\nolinkurl{#1}}

\paragraph{Scope (formal mathematics, not phenomenology).}
This note records machine-checked regional closed forms, projection identities,
and logical equivalences in the public HQIV Lean library.  It does \emph{not}
assert phenomenological predictions, a scale witness, or patch-closed
completeness in the TUFT+SM sense.  When classical notation appears ($\Gamma$,
sin/cos, heat flow, de~Bruijn--Newman families), it serves as a
\emph{comparison dictionary} for readers trained in analytic number theory; the
proved objects are discrete harmonic sums, shell weights, prime Euler products,
and explicit predicates on the open strip.  No fundamental continuum is required
for the theorems cited here; a literal continuum limit on the load-bearing shell
index would erase the harmonic--$\Delta$ closure that motivates the construction
(\cite{hqiv-so8-paper}).


\providecommand{\leanfile}[1]{\nolinkurl{#1}}
\providecommand{\leanid}[1]{\nolinkurl{#1}}

\paragraph{Scope (formal mathematics, not phenomenology).}
This note records machine-checked regional closed forms, projection identities,
and logical equivalences in the public HQIV Lean library.  It does \emph{not}
assert phenomenological predictions, a scale witness, or patch-closed
completeness in the TUFT+SM sense.  When classical notation appears ($\Gamma$,
sin/cos, heat flow, de~Bruijn--Newman families), it serves as a
\emph{comparison dictionary} for readers trained in analytic number theory; the
proved objects are discrete harmonic sums, shell weights, prime Euler products,
and explicit predicates on the open strip.  No fundamental continuum is required
for the theorems cited here; a literal continuum limit on the load-bearing shell
index would erase the harmonic--$\Delta$ closure that motivates the construction
(\cite{hqiv-so8-paper}).


\providecommand{\leanfile}[1]{\nolinkurl{#1}}
\providecommand{\leanid}[1]{\nolinkurl{#1}}

\paragraph{Scope (formal mathematics, not phenomenology).}
This note records machine-checked regional closed forms, projection identities,
and logical equivalences in the public HQIV Lean library.  It does \emph{not}
assert phenomenological predictions, a scale witness, or patch-closed
completeness in the TUFT+SM sense.  When classical notation appears ($\Gamma$,
sin/cos, heat flow, de~Bruijn--Newman families), it serves as a
\emph{comparison dictionary} for readers trained in analytic number theory; the
proved objects are discrete harmonic sums, shell weights, prime Euler products,
and explicit predicates on the open strip.  No fundamental continuum is required
for the theorems cited here; a literal continuum limit on the load-bearing shell
index would erase the harmonic--$\Delta$ closure that motivates the construction
(\cite{hqiv-so8-paper}).


\section{Introduction}\label{sec:intro}

The Riemann zeta function is usually introduced as a Dirichlet series and then
analytically continued.  On the open critical strip $0<\Re s<1$, the continued
formula is the standard sin--cos--$\Gamma$ functional-equation assembly:
products of $\Gamma$, $\pi$, $\sin(\pi s/2)$, and $\cos(\pi s/2)$ that look
over-determined and singular until trivial zeros and continuation slots are
separated (\S\ref{sec:real-line}).  This note gives a geometric interpretation of
that assembly: the functional equation is reflection symmetry
$\sigma\leftrightarrow 1-\sigma$ rotated by $45^\circ$; a fixed diagonal
component separates from a variable \emph{equator factor}
$(2\sigma-1)/\sqrt2$; and the divergent harmonic series of the closure paper
forces a phase-lift connector $\Delta$ that closes $\mathfrak{so}(3)+\Delta$ to
$\mathfrak{so}(4)$ on a four-dimensional quaternion carrier.

This note records what is already proved in Lean:
\begin{enumerate}[label=(\roman*),nosep]
\item the harmonic--$\Delta$--SO(4) closure bridge;
\item the harmonic $\Rightarrow$ prime Euler product $\Rightarrow$ $\zeta$
      $\Rightarrow$ $\Lambda=-\zeta'/\zeta$ pipeline for $\Re s>1$;
\item regional closed forms for $\zeta(s)$ matching Mathlib;
\item sine/cosine factor cancellation at even $\pi$-multiples and rotation ratios
      on odd / fractional arguments;
\item exclusion lemmas ruling out nontrivial zeros in half-planes and at
      negative integer slots;
\item a factorization $\zeta(s)=h(s)\cdot(2\sigma-1)/\sqrt2$ on the open strip
      off $\sigma=\tfrac12$, with $h=\texttt{interiorStripH}$;
\item dependence on $\Re s$ only: the equator factor is a function of
      $\sigma=\Re s$ and cannot globally identify $\zeta$ with the projection
      formula (\S\ref{sec:sigma-scope});
\item FE + Schwarz quadruplet orbits for $\zeta$-zeros, with unconditional FE
      pairs and conditional full quadruplets (\S\ref{sec:quadruplet});
\item the shared $\tfrac12$ slope bridge: exact equivalence of RH $\wedge$
      Goldbach with a single geometric hypothesis (\S\ref{sec:rh-goldbach});
\item the discrete/continuum heat-flow comparison
      $c\lVert H_t\rVert\le\lVert\texttt{HQIVDeformedSum}\rVert$ with exact
      phase alignment (\S\ref{sec:dbn-comparison}), and the de Bruijn--Newman
      constant ported as a concrete Lean definition with Rodgers--Tao and
      Newman/dBN as named statements (\S\ref{sec:dbn-port});
\item the tangent law of the quadruplet orbit: FE inverts and Schwarz
      conjugates $T(s)=\tan(\pi s/2)$, and the critical line is exactly the
      unimodular-tangent locus (\S\ref{sec:tangent-orbit});
\item zeros as spectral cancellations of the Euler product: no cancellation
      in the product region, breakdown-locus confinement, and resonance
      blow-up of $\lVert1/\zeta\rVert$ at every zero
      (\S\ref{sec:spectral-cancellation});
\item resonance channeling: the two-channel dichotomy at every zero, the
      equivalence \texttt{InteriorAssemblyNonzero}\,$\Leftrightarrow$\,RH,
      the SO(4) spectral-line transformer $n\mapsto n^{-s}$ with FE pairing
      $\lVert n^{-s}\rVert\lVert n^{-(1-s)}\rVert=1/n$, and
      RH $\Leftrightarrow$ square-root spectral weights at every zero
      (\S\ref{sec:resonance-channeling});
\item vertical axis rigidity: composing the FE point-reflection with any
      second mirror axis yields a translation that marches zeros out of the
      strip, so the zero set admits at most one vertical axis and it must be
      $\Re s=\tfrac12$ (\S\ref{sec:axis-rigidity});
\item same-height orbit collapse: the quadruplet partner $1-\overline{\rho}$
      shares $\rho$'s height and merges with it exactly on the line, giving
      RH $\Leftrightarrow$ ``one zero per height'' --- unconditionally in
      the $\lvert\Im\rvert$ form, and with explicit Schwarz hypothesis in
      the strict form (\S\ref{sec:same-height});
\item polar projection collapse: the polar separation
      $1-\overline{s}-s=-\sqrt2\cdot\texttt{so4CriticalFactor}(s)$ is the
      equator factor itself, and single-point polar projection on the zero
      set is RH (\S\ref{sec:polar-collapse});
\item the zero--holonomy--Goldbach chain: every zero activates every
      midpoint pair, each pair carries Hopf holonomy support, and the
      half-slope bridge rewrites as polar collapse $\wedge$ parity
      (\S\ref{sec:zero-holonomy-chain});
\item the TUFT nested frame tower: level-$N$ spectral frames on odd spheres
      $S^{2N-1}$, the harmonic radius law
      $\texttt{frameNormSq}=H_N\Leftrightarrow\Re s=\tfrac12$ at every
      level, divergence of the radius ladder on the line, and RH
      $\Leftrightarrow$ harmonic radii at every zero
      (\S\ref{sec:tuft-tower});
\item the Diophantine transformer: the level-$N$ diagonal operator
      $D_N(s)=\mathrm{diag}(1^{-s},\dots,N^{-s})$ whose trace ladder
      converges to $\zeta$, whose Hilbert--Schmidt norm is the tower
      radius, whose FE composition is the fixed rational harmonic operator
      $\mathrm{diag}(1/n)$, and whose adjoint law
      $D_N(1-s)=D_N(s)^{\dagger}\Leftrightarrow\Re s=\tfrac12$ gives RH
      $\Leftrightarrow$ self-adjointness across FE at every zero
      (\S\ref{sec:diophantine-transformer});
\item the octonionic associator channel: the transformer family is
      provably commutative and normal (the formal reason the diagonal
      attack saturates), while the non-associative carrier escapes —
      the triple $(e_1,e_2,e_4)$ has $[e_1,e_2,e_4]=-e_3-e_4\neq0$ on the
      frozen Fano tables, prime spectral weights ride the associator
      trilinearly, the channel $2(abc)^{-2\sigma}$ never vanishes
      (zeros kill the abelian assembly, never the torsion), and RH
      $\Leftrightarrow$ square-root torsion weight at every zero
      (\S\ref{sec:associator-channel});
\item S$^7$ torsion cancellation under SO(8): the torsion vector is a
      positive multiple of one fixed carrier direction (the S$^7$ shadow
      never moves), the FE defect is odd with a chirality flip across the
      line, FE partners' torsions share an SO(8) orbit iff they are
      equal iff $\Re s=\tfrac12$, and RH $\Leftrightarrow$ the torsion
      defect vanishes at every zero
      (\S\ref{sec:s7-torsion});
\item the critical-line carrier bundle: the strip embeds as a cylinder
      $\sigma\times S^1$, but only at $\Re s=\tfrac12$ is the lift the full
      Hopf circle ($r=1$, $f=0$, rolling map); off the line the fiber shrinks
      ($r<1$), the tangent sorts the half-planes hyperbolically, and
      quadruplet collapse fails---whole-strip classification is conditional
      (\S\ref{sec:analytic-strip-lift});
\item the log edge: polar decoupling
      $n^{-s}=n^{-\sigma}\cdot e^{-it\log n}$ (modulus knows only
      $\sigma$, phase only $t$), Diophantine independence of prime
      phase speeds ($k\log p=m\log q\Rightarrow k=m=0$, from unique
      factorization), two prime phases pin the height exactly, the
      Goldbach circle $pq+(q-N)^2=N^2$ as literal additive curvature
      with phase-speed cap $\log(pq)\le2\log N$, and — under RH — two
      prime phases as a complete address for a zero
      (\S\ref{sec:log-edge}).
\end{enumerate}
The remaining RH step is \emph{not} continuation; it is proving that the assembly
$h$ does not vanish at nontrivial zeros off the line---equivalently, upgrading
orbit cancellation (functional-equation pairs) to pointwise cancellation at
actual $\zeta$ zeros.  Goldbach is the same geometric half-slope on the additive
prime channel, equivalent only under the holonomy bridge hypothesis.

\subsection{Why zeta functions appear in HEP and other HQIV integrations}
\label{sec:hep-zeta-bridge}

HQIV does not use ``zeta'' as a single object.  Three layers recur across the
library and paper series; this note isolates the \emph{classical projection}
layer and explains why the same vocabulary appears in high-energy phenomenology
readouts without implying that every ``zeta'' is Riemann's function.

\paragraph{Sector zeta on the patch carrier.}
In the TUFT+SM synthesis \cite{hqiv-tuft-sm-lagrangian-paper}, \emph{sector
zeta} means fibre-torsion subleading spectral bodies on nested Hopf shells:
Beltrami-sector determinants, \texttt{zetaHQIVTerm} phase on the null-lattice
carrier, and holonomy discharge bundles on $\mathfrak{so}(8)$.  These are
\emph{finite} shell-indexed readouts tied to mass rungs and coupling ladders on
the accessible patch---not the complex-analytic $\zeta(s)$ studied here.

\paragraph{Classical $\zeta(s)$ as projection geometry (this note).}
The present manuscript shows that Mathlib's Riemann zeta is already a
\emph{regional closed form} of an SO(4) projection: harmonic shell weights sum
to $\zeta(s)$ for $\Re s>1$; prime Euler products and $\Lambda=-\zeta'/\zeta$
follow on the same region; the open strip is a sin--cos--$\Gamma$ assembly with
a $45^\circ$ equator factor $(2\sigma-1)/\sqrt2$.  The proved polar law
$n^{-s}=n^{-\sigma}\cdot e^{-it\log n}$ (\S\ref{sec:log-edge}) is the
structural reason modulus-based equivalent conditions throughout this note
reformulate RH
without importing phase data into $\sigma$.

\paragraph{HEP multichannel and decay formulas.}
The heavy-flavour decay programme \cite{hqiv-hep-decay-paper} uses the same
\emph{spectral--log} pattern at the formula layer: discrete multichannel
branching rules, three-ledger $\beta$ tipping, and rational EM contact
coefficients are computed from composite traces and shell weights without
partial-width inputs.  Zeta language enters because \textbf{(i)}~prime and
shell indices carry multiplicative weights $n^{-s}$ (or their real-axis
cousins) in every channel; \textbf{(ii)}~logarithmic phase speeds $\log p$,
$\log N$ tie additive midpoint data (Goldbach circle, \S\ref{sec:log-edge}) to
multiplicative pooling (quarkonium channels, joint products $pq$); and
\textbf{(iii)}~continuation and resonance (zeros of assembled factors, blow-up
of $1/\zeta$, spectral cancellation profile, \S\ref{sec:spectral-cancellation})
is the same geometric template as ``channel factors vanish where the global
assembly vanishes.''  This note does not reproduce HEP numerics; it supplies
the proved analytic spine that makes the shared vocabulary precise: zeta
functions in HQIV integrations are projection formulas for prime/shell spectral
data, not ad hoc naming.

\paragraph{Continuum notation elsewhere.}
Heat-flow, de~Bruijn--Newman, and O-Maxwell continuum charts in sibling
manuscripts are \emph{comparison dictionaries} for classical literature---as
in \texttt{readout\_dictionary\_messaging.tex} above, not load-bearing
continuum extensions.  The discrete harmonic ladder and shell index remain the
proved ontology; a literal continuum limit on those indices would erase the
predictions (\cite{hqiv-so8-paper}).

\section{Harmonic divergence forces $\Delta$ and SO(4)}\label{sec:harmonic-delta}

The closure spine \cite{hqiv-so8-paper} proves a uniform lower bound
\[
  H_n := \sum_{m=0}^{n-1}\frac{1}{m+1}
  \;\le\;
  K(n) := \sum_{m=0}^{n-1}\underbrace{\frac{1}{m+1}\bigl(1+\alpha\log(m+1)\bigr)}_{\text{shell density}},
  \qquad \alpha=\tfrac35,
\]
with $K$ strictly increasing (\leanfile{Hqiv/Story/S3ClosureDeltaLiftBridge.lean},
\leanfile{Hqiv/Geometry/OctonionicLightCone.lean}).
The decomposition
\[
  K(n)=H_n+\alpha\sum_{m=0}^{n-1}\frac{\log(m+1)}{m+1}
\]
is the exact curvature-integral split
(\leanid{curvature\_integral\_eq\_harmonic\_plus\_alpha\_log}).

$\Delta$ is \emph{not} identical to $H_n$, but it is forced because the
normalized channel $\Omega=K/K_*$ diverges.  On the toy four-dimensional
carrier the Lie span of $\mathfrak{so}(3)$ together with the connector
$\Delta_4=J_{14}$ generates full $\mathfrak{so}(4)$
(\leanid{so3\_delta\_lifts\_to\_so4}).

\subsection{Third harmonic orbit and the $1+\alpha/3$ projection}

The third partial sum $H_3=1+\tfrac12+\tfrac13$ is the first shell carrying an
explicit $1/3$ harmonic term.  Against the HQIV imprint $\alpha=3/5$ the
projection multiplier
\[
  1+\frac{\alpha}{3}=1+\frac15=\frac65=1.2
\]
is proved exactly (\leanfile{Hqiv/Story/S3HarmonicDeltaEvenOrbit.lean},
\leanid{so4HarmonicThirdProjection\_eq\_one\_point\_two}).
On the eight-dimensional shell the unit $7$-sphere area proxy is $\pi^4/3$;
halving at the $45^\circ$ equator gives $\pi^4/6$
(\leanid{so4EquatorHalfArea\_eq\_pi\_four\_sixths}), parallel to the even
Bernoulli slot $\zeta(2)=\pi^2/6$ at $\pi^2$ power.

\section{Forty-five degree projection}\label{sec:45-proj}

\begin{definition}[Functional-equation pair and $45^\circ$ slots]
For $\sigma\in\mathbb{R}$ define the pair $(\sigma,1-\sigma)$ and the
rotated coordinates
\begin{align*}
  \texttt{rot45Diag}(\sigma,1-\sigma) &= \frac{\sigma+(1-\sigma)}{\sqrt2}=\frac1{\sqrt2},\\
  \texttt{rot45Free}(\sigma,1-\sigma) &= \frac{\sigma-(1-\sigma)}{\sqrt2}=\frac{2\sigma-1}{\sqrt2}.
\end{align*}
At $\theta=\pi/4$ one has $\cos\theta=\sin\theta=\sqrt2/2$ (not $\pi$); $\pi$
enters through the functional-equation angle $\pi s/2$ below.
\end{definition}

\begin{theorem}[Critical line as equator]\label{thm:critical-line}
The free coordinate vanishes iff $\sigma=\tfrac12$:
\[
  \texttt{rot45Free}(\sigma,1-\sigma)=0
  \quad\Longleftrightarrow\quad
  \sigma=\tfrac12.
\]
A nonzero rotation offset $\pi/4+\delta$ with $\sin\delta\neq0$ moves the
vanishing locus off $\sigma=\tfrac12$
(\leanfile{Hqiv/Story/S3RotationRigidity.lean}).
\end{theorem}

The complex lift of the equator factor is the \textbf{SO(4) critical factor}
\[
  \texttt{so4CriticalFactor}(s)=\frac{2\sigma-1}{\sqrt2},
  \qquad s=\sigma+i t,
\]
which vanishes exactly on $\Re s=\tfrac12$
(\leanfile{Hqiv/Story/S3ZetaClosedForm.lean}).

\subsection{The equator factor depends only on $\Re s$}\label{sec:sigma-scope}

The exact $45^\circ$ projection is distinguished: other rotation angles move
the vanishing locus off $\sigma=\tfrac12$
(\leanfile{Hqiv/Story/S3RotationRigidity.lean}).  What the projection formula
\emph{cannot} decide is also proved:
\begin{itemize}[nosep]
\item \textbf{$\sigma$-blindness.} \texttt{exactTwiddleReadout} and
      \texttt{so4CriticalFactor} depend only on $\sigma=\Re s$, not on height
      $t=\Im s$
      (\leanfile{Hqiv/Story/S3SigmaReadoutScope.lean},
      \leanid{sigmaReadout\_blind\_to\_imaginary\_height}).
\item \textbf{Identification $\not\Rightarrow$ constraint.}  Pointwise equality
      $\zeta(s)=\texttt{exactTwiddleReadout}(s)$ transfers zero iff at $s$, but
      does not force $\sigma=\tfrac12$ without that identification
      (\leanid{zeta\_eq\_readout\_same\_zero\_set}).
\item \textbf{No global identification.}  $\zeta(s)=\texttt{exactTwiddleReadout}(s)$
      for all $s$ fails: $\zeta(0)=-\tfrac12$ but the projection formula gives $-1/\sqrt2$
      (\leanid{not\_globally\_zeta\_eq\_real\_twiddle\_readout}).
\item \textbf{Affine form.}  Off the even lift,
      $\texttt{so4CriticalFactor}(z)=(2\Re z-1)/\sqrt2$
      (\leanid{so4CriticalFactor\_eq\_affine\_re}); this is real-valued in $\sigma$
      and not holomorphic in $s$.
\end{itemize}
Proving regional $\zeta$ closed forms therefore \emph{relocates} zeros onto the
construction without transferring the equator vanishing locus onto $\zeta$.

\subsection{Mirrored twiddle pairs}\label{sec:mirrored-twiddle}

The functional equation supplies reflection $\sigma\leftrightarrow1-\sigma$.
Rather than treating twiddle angles as independent knobs, one may pair angles
that are mirrors under the same reflection:
\[
  \theta\;\longleftrightarrow\;\frac{\pi}{2}-\theta.
\]
\leanfile{Hqiv/Story/S3MirroredTwiddlePair.lean} formalizes two complementary
readouts.

\paragraph{Rotation twiddle (paper default).}
On the FE pair $(\sigma,1-\sigma)$ the free coordinate after rotation by
$\theta$ is
\[
  \texttt{rotFree}_\theta(\sigma)=\sigma(\cos\theta+\sin\theta)-\sin\theta,
\]
with vanishing locus the \textbf{projection line}
$\sigma_0(\theta)=\sin\theta/(\cos\theta+\sin\theta)$
(\leanfile{Hqiv/Story/S3RotationRigidity.lean}).
Its mirror at $\pi/2-\theta$ has zero locus
\[
  \sigma_0\!\left(\frac{\pi}{2}-\theta\right)=1-\sigma_0(\theta),
\]
so the two candidate axes are symmetric about $\sigma=\tfrac12$
(\leanid{projectionLine\_mirror}).
Across the FE orbit the pair cancels pairwise:
\[
  \texttt{rotFree}_\theta(\sigma,1-\sigma)
  +\texttt{rotFree}_{\pi/2-\theta}(1-\sigma,\sigma)=0
\]
(\leanid{rotFree\_mirror\_pair\_sum}); on a fixed pair $(\sigma,1-\sigma)$ the
sum and difference separate into
\[
  (\cos\theta+\sin\theta)(2\sigma-1)
  \quad\text{and}\quad
  \cos\theta-\sin\theta,
\]
so the critical line is the \textbf{coincidence locus} where the two channels
agree in the $\sigma$-only sense: the difference is constant in $\sigma$ and
vanishes only at $\theta=\pi/4$, when both angles coincide and
$\sigma_0(\pi/4)=\tfrac12$.

\paragraph{Affine twiddle (weighted $\sigma$ form).}
The weighted combination
\[
  \texttt{affineTwiddle}_\theta(\sigma)=\cos\theta\,\sigma+\sin\theta\,(1-\sigma)
\]
is the mirror swap
$\sin\theta\,\sigma+\cos\theta\,(1-\sigma)$ at $\pi/2-\theta$
(\leanid{affineTwiddle\_mirror\_angle}).
When $\sin\theta\neq\cos\theta$ its zero is
$\sigma=\sin\theta/(\sin\theta-\cos\theta)$, again mirrored by
$\sigma\mapsto1-\sigma$ (\leanid{affineTwiddleZero\_mirror}).
At $\theta=\pi/4$ this family is \emph{degenerate}: $\cos\theta=\sin\theta$ so
the readout is the constant $\sin(\pi/4)\neq0$ and never vanishes
(\leanid{affineTwiddle\_pi\_div\_four\_ne\_zero}); the $45^\circ$ critical-line
pin is native to the rotation/normalized form, not this affine parameterization.

\paragraph{Example: $30^\circ$ and $60^\circ$.}
$\pi/6$ and $\pi/3$ are a mirror pair.
Their affine zero loci are far from $\sigma=\tfrac12$ (e.g.\ at $\pi/6$ the zero
is $(1/2)/(1/2-\sqrt3/2)$), but they satisfy the exact swap
$\sigma_0(\pi/3)=1-\sigma_0(\pi/6)$
(\leanid{affineTwiddleZero\_pi\_six\_pi\_three\_mirror}).
Near-$45^\circ$ rotation angles instead give projection lines
$\sigma=\tfrac12\pm\varepsilon$ --- the regime where both mirrors hug the
critical line.

\paragraph{Interaction with sine/cosine slots and vertical rigidity.}
The FE channel already carries the angle $\pi s/2$ in
$\sin(\pi s/2)$ and $\cos(\pi s/2)$ (\S\ref{sec:quaternion}); mirrored twiddles
add a second angular degree on the real $\sigma$ pair, enriching the
even/odd/fractional projection trichotomy without replacing it.
Each individual twiddle still defines a straight vertical candidate axis; if
two such mirrors are imposed together with FE closure, the same translation
argument as in \S\ref{sec:axis-rigidity} forces any common strip zero locus onto
the unique fixed line $\Re s=\tfrac12$.
Off the line the pair distinguishes \emph{orbit} cancellation (FE partners) from
\emph{pointwise} coincidence on $\sigma=\tfrac12$ --- a two-channel sharpening
of the orbit-vs-pointwise gap (\S\ref{sec:orbit-gap}).

\subsection{Higher twiddles as probes of the factorization engine}\label{sec:higher-twiddle-probes}

The linear-twiddle family and its cubic specializations are \textbf{not} intended
to yield new information about the location of zeros on their own loci.  Their
purpose is to illuminate the \textbf{factorization engine} --- the decomposition
of $\zeta(s)$ into an interior assembly times a geometric readout factor --- when
the readout is taken from channels other than the distinguished $45^\circ$
projection.

\paragraph{Main channel (paper default).}
Off the critical line on the open strip,
\[
  \zeta(s)=h(s)\cdot\texttt{so4CriticalFactor}(s)
  =h(s)\cdot\frac{2\sigma-1}{\sqrt2}
\]
(\S\ref{sec:explicit-h}, \eqref{eq:h-explicit}): the interior assembly
$h=\texttt{interiorStripH}$ absorbs everything that is not the equator readout.

\paragraph{Higher-twiddle readouts.}
For a coefficient $c\in(0,1)$ the linear free readout is
\[
  \text{free}_c(\sigma)=c\cdot\sigma+(1-c)\cdot(1-\sigma),
\]
with zero on the vertical line $\sigma_c=(c-1)/(2c-1)$ when $2c-1\neq0$.
Functional-equation reflection interchanges $c$ with $1-c$, producing mirrored
pairs (\S\ref{sec:mirrored-twiddle}).
Specializing to the plastic constant via $c_\rho=1-\rho^{-1}$ (unique $\rho>1$ with
$\rho^3=\rho+1$) produces the normalized cubic twiddle
\texttt{so4PlasticCubicFree} whose zero locus is the explicit algebraic number
\[
  \sigma_\rho=\frac{1}{2-\rho}
\]
(\leanfile{Hqiv/Story/S3CubicPlasticTwiddle.lean}:
mirror symmetry, normalization, and FE partner
\leanid{so4PlasticCubicFree\_one\_sub};
patch factorization off the probe locus in
\leanfile{Hqiv/Story/S3HigherTwiddleFactorizationProbe.lean}).
This locus lies far from $\Re s=\tfrac12$; the construction is a \emph{controlled
perturbation} of the readout geometry, not a competing critical-line candidate.

\paragraph{Multi-channel laboratory.}
Because cubic (and, more generally, higher-degree) twiddles are not expected to
constrain zeros directly on their own axes, they serve as probes of
\textbf{cross-channel} behavior inside the factorization architecture:
\begin{itemize}[nosep]
\item How does the interior assembly change when the geometric factor is
      $\texttt{so4PlasticCubicFree}$ (or another $\text{free}_c$) instead of
      \texttt{so4CriticalFactor}?
\item What is the coupling between the main $45^\circ$ channel and a higher-twiddle
      channel when both are measured at the same strip point?
\item Can one define a controlled cross-talk term supported only on semiprime
      spectral weights?
\end{itemize}
The patch-level target, not yet formalized, is a family
\[
  \zeta(s)=h_c(s)\cdot\text{free}_c(\Re s)
  \quad\text{(off the zero locus of channel $c$)},
\]
parallel to \eqref{eq:h-explicit} but with readout channel $c$ replacing the
$45^\circ$ equator factor.

\paragraph{Semiprime-mediated influence on the critical line.}
Any influence a higher-twiddle channel exerts on points lying on the $45^\circ$
critical line is expected to be \textbf{semiprime mediated}: it can reach the
main channel only through spectral weights supported on integers with exactly
two prime factors (semiprimes $pq$ with $p\neq q$ in the square-free case).
This mirrors the arithmetic difficulty of locating or controlling semiprimes in
the additive Goldbach channel (\S\ref{sec:rh-goldbach}): midpoint pairs
$p+q=2N$ are easy to \emph{certify} once found, while products $pq$ are the
composite objects that factorization-side bookkeeping must resolve
(\leanfile{Hqiv/Geometry/GoldbachG2Parity.lean}, composite-branch and
\texttt{G$_2$}-triple decomposition lemmas).
Consequently higher twiddles may affect the factorization engine at
critical-line points only \emph{indirectly and weakly} --- not as new zero
locators, but as cross-talk that must pass through the hardest multiplicative
spectral lines.

\paragraph{Open technical tasks.}
\begin{enumerate}[label=(\roman*),nosep]
\item Define $h_c(s)$ (or $h_\rho$) relative to a higher-twiddle readout, matching
      the interior assembly role of \texttt{interiorStripH}.
\item Prove the corresponding patch-level factorization inside finite spectral
      frames of level $N$.
\item Make precise the semiprime support of any cross-talk between the $45^\circ$
      channel and higher-twiddle channels (a leakage term vanishing on primes and
      higher-composite weights).  Lean already types finite-frame leakage with
      \texttt{SemiprimeSupportedWeight} and proves prime spectral lines carry
      zero weight (\leanfile{Hqiv/Story/S3HigherTwiddleFactorizationProbe.lean},
      \leanid{semiprime\_weight\_kills\_prime\_spectral\_line}).
\end{enumerate}

\section{Quaternion axis projection}\label{sec:quaternion}

On the unit imaginary quaternion shell $S^3\subset\mathbb{H}$ the six axis poles
$\pm i,\pm j,\pm k$ project to $\pm1/\sqrt2$ on single axes and cancel as
antipodal orbit pairs (\leanfile{Hqiv/Story/S3SixPolesResidual.lean},
\leanfile{Hqiv/Story/S3PrimeAxisCancellation.lean}).

\begin{definition}[Sine and cosine factors in the FE channel]
Define the sine and cosine factors
\[
  \mathcal{S}(s)=\sin\frac{\pi s}{2},
  \qquad
  \mathcal{C}(s)=\cos\frac{\pi s}{2}.
\]
These are the oscillating factors in the functional-equation odd channel;
$\mathcal{C}$ is \texttt{zetaSinCosFactor} in Lean.
\end{definition}

\begin{theorem}[Even $\pi$-multiples: sine vanishes]\label{thm:even-cancel}
For $k\in\mathbb{N}$ at $s=2k$:
\[
  \mathcal{S}(2k)=\sin(k\pi)=0,
  \qquad
  \mathcal{C}(2k)=\cos(k\pi)=(-1)^k.
\]
\end{theorem}

\begin{theorem}[Odd integers: cosine cancels]\label{thm:odd-pure-sin}
For $k\in\mathbb{N}$ at $s=2k+1$:
\[
  \mathcal{C}(2k+1)=0,
  \qquad
  \mathcal{S}(2k+1)=(-1)^k.
\]
\end{theorem}

\begin{theorem}[Fractional strip: rotation ratio]\label{thm:frac-ratio}
For $0<\Re s<1$ with $\mathcal{C}(s)\neq0$,
\[
  \frac{\mathcal{S}(s)}{\mathcal{C}(s)}=\tan\frac{\pi s}{2}.
\]
At $\sigma=\tfrac12$ both slots equal $\sqrt2/2$ and the ratio is $1$
(\leanid{zetaAxisRotationRatio\_at\_half}).
\end{theorem}

Theorems~\ref{thm:even-cancel}--\ref{thm:frac-ratio} are proved in
\leanfile{Hqiv/Story/S3ZetaAxisRotationProjection.lean}.
They explain why \textbf{even} $\zeta$-values admit $\pi$-sector closed forms
while \textbf{odd} positive integers (e.g.\ $\zeta(3)$) occupy the Ap\'ery /
plastic continuation slot with no elementary closed form in Mathlib.

\section{Analytic strip lift: cylinder $\sigma\times S^1$ in $\mathbb{H}$}
\label{sec:analytic-strip-lift}

The critical strip $0<\sigma<1$ embeds as a \emph{cylinder cover}: at fixed
$\sigma$ the height $t=\Im s$ rolls a Hopf fiber on $S^3\subset\mathbb{H}$,
while $\sigma$ enters through the $45^\circ$ free coordinate of
\S\ref{sec:45-proj}.
The geometric punch is not ``the whole strip is a flat cylinder'' but
\textbf{what is special about the critical line}: there the free coordinate
vanishes, the fiber radius is maximal, the lift is the rolling map, and the
tangent readout is unimodular with quadruplet collapse; off the line the
analytic slots carry hyperbolic ($\cosh/\sinh$) sorting and the Hopf circle is
strictly shrunk.
Lean packages the contrast in
\leanfile{Hqiv/Story/S3CriticalLineCarrierBundle.lean}; the lift and
conditional whole-strip classification remain in
\leanfile{Hqiv/Story/S3AnalyticStripLift.lean} and
\leanfile{Hqiv/Story/S3AnalyticStripClassification.lean}.

\subsection{Lift and Hopf fiber}

\begin{definition}[Analytic strip lift]
For $\sigma,t\in\mathbb{R}$ define the \textbf{free coordinate}
\[
  f(\sigma)=\frac{2\sigma-1}{\sqrt2}=\texttt{rot45Free}(\sigma,1-\sigma),
\]
the \textbf{Hopf fiber radius}
$r(\sigma)=\sqrt{1-f(\sigma)^2}$, and the lift
\[
  \texttt{stripAnalyticLift}(\sigma,t)
  = \bigl(0,\; f(\sigma),\; r(\sigma)\cos t,\; r(\sigma)\sin t\bigr)
\]
on the $p_0=0$ slice of $S^3$.
For $s=\sigma+it$ set $\texttt{stripPointLift}(s)=\texttt{stripAnalyticLift}(\sigma,t)$.
\end{definition}

\begin{theorem}[On-shell and $\sigma$-equator]\label{thm:lift-on-s3}
For $0\le\sigma\le1$ and all $t$, the lift lies on $S^3$
(\leanid{strip\_analytic\_lift\_on\_s3}).
The free coordinate vanishes exactly on the critical line:
$f(\sigma)=0\Leftrightarrow\sigma=\tfrac12$
(\leanid{strip\_sigma\_free\_coord\_vanishes\_iff}).
At $\sigma=\tfrac12$ the lift agrees with the rolling map of
\leanfile{Hqiv/Story/S3StripRollingProjection.lean}:
$\texttt{stripAnalyticLift}(\tfrac12,t)=\texttt{stripRollingMap}(t)$
(\leanid{strip\_analytic\_lift\_eq\_rolling}).
\end{theorem}

\subsection{Critical line carrier bundle}\label{subsec:critical-line-carrier}

\begin{theorem}[Line is the unique maximal-circle, unimodular locus]\label{thm:line-carrier-chain}
On the open strip $0<\sigma<1$, the following are equivalent
(\leanid{critical\_line\_carrier\_iff\_chain}):
\[
  \Re s=\tfrac12
  \;\Longleftrightarrow\;
  f(\sigma)=0 \;\wedge\; r(\sigma)=1
  \;\wedge\; \texttt{criticalLineDeviation}(s)=0
  \;\wedge\; \texttt{so4CriticalFactor}(s)=0
  \;\wedge\; \lVert\tan(\pi s/2)\rVert=1
  \;\wedge\; T(s)^{-1}=\overline{T(s)}.
\]
Here $f(\sigma)=\texttt{stripSigmaFreeCoord}(\sigma)$ and
$r(\sigma)=\texttt{stripHopfFiberRadius}(\sigma)$.
The lift on the line is the rolling map
(\leanid{critical\_line\_lift\_is\_rolling}).
\end{theorem}

\begin{theorem}[Off-line carrier tax]\label{thm:off-line-carrier}
For $0\le\sigma\le1$ with $\sigma\neq\tfrac12$,
$f(\sigma)\neq0$ and $r(\sigma)<1$
(\leanid{off\_line\_carrier\_shrinks\_hopf\_fiber},
\leanid{strip\_hopf\_fiber\_radius\_lt\_one\_iff}).
On the open strip with $\Re s\neq\tfrac12$, the tangent is not unimodular,
the quadruplet does not collapse, and $\sigma$ sorts into the disk or
exterior of the unit circle
(\leanid{off\_line\_hyperbolic\_tangent\_sorts},
\leanid{critical\_strip\_off\_line\_quadruplet\_does\_not\_collapse}).
Cancellation off the line requires the full offset
$f(\sigma)+r(\sigma)(\cos t+\sin t)=0$, not merely $\cos t+\sin t=0$
(\leanid{strip\_analytic\_cancellation\_off\_line\_needs\_free}).
\end{theorem}

\begin{theorem}[Path C on the line via rolling]\label{thm:line-pathC-rolling}
Under \texttt{StripAnalyticLiftMatches} and \texttt{ZetaEqualsS3ResidualAt},
for $\Re s=\tfrac12$,
\[
  \zeta(s)=0
  \;\Longleftrightarrow\;
  \texttt{ZeroProducingOrbit}(\texttt{stripRollingMap}(t))
  \;\wedge\;
  \texttt{PathEStripHolonomyCloses}(s)
\]
(\leanid{critical\_line\_pathC\_classification\_via\_rolling});
equivalently, only the rolling orbit is needed on the line
(\leanid{critical\_line\_pathC\_rolling\_only}).
\end{theorem}

The hyperbolic character off-line is in the \emph{analytic} sin/cos slots
($\cosh/\sinh$ factors cancel only in
$\lVert\sin\rVert^2-\lVert\cos\rVert^2=-\cos(\pi\sigma)$,
\leanid{normSq\_proj\_sin\_sub\_cos}), not in the constraint
$f(\sigma)^2+r(\sigma)^2=1$ (a unit circle in the lift chart,
\leanid{strip\_lift\_free\_radius\_unit\_circle}).

\subsection{$\mathbb{C}$ cover versus $\mathbb{H}$ circle}

In $\mathbb{C}$ the height $t$ is a \emph{non-compact} cover coordinate.
On the lift, $t$ and $t+2\pi n$ define the same point on $S^3$
(\leanid{strip\_analytic\_lift\_period},
\leanid{sameStripHeight}).
The Hopf-base trace $(r(\sigma)\cos t,r(\sigma)\sin t)$ is likewise
$2\pi$-periodic in $t$ (\leanid{strip\_hopf\_circle\_period}).
Thus each fixed $\sigma$ carries an $S^1$ fiber inside $\mathbb{H}$ while
$\mathbb{C}$ records the universal cover --- the geometric form of
``countably infinite heights, compact phase carrier'' discussed in the
spectral literature without the quaternion chart.

\subsection{Harmonic-shell arc sweep (not heights $t$)}\label{subsec:harmonic-shell-counting}

Counting by scanning $t\in\mathbb{R}$ conflates the universal cover with a counter.
On the line, $t\neq t+2\pi$ yet the same fiber class is recovered
(\leanid{strip\_height\_not\_injective\_on\_circle},
\leanid{strip\_rolling\_map\_respects\_height\_cover}).

The promoted counting axis pairs harmonic depth $H_n$ with angular \emph{arc
partitions}.  At shell $n>0$ the Hopf circle is cut into $n$ arcs of width
$2\pi/n$, indexed by $k$ at polar angle $2\pi k/n$
(\leanid{shellSweepAngle}, \leanid{shellArcWidth}).
Finer shells refine the sweep when $n\mid m$
(\leanid{sweep\_refines\_coarser\_angle}, \leanid{sweep\_arc\_width\_refines}).
Slots form a countable union of finite partitions:
$\texttt{HarmonicShellSlot}=\Sigma_n\,\texttt{Fin}\,n$ injects into $\mathbb{N}$
(\leanid{encodeHarmonicShellSlot\_injective}).

Monotone ticks $1,2,3,\ldots$ order shells by cumulative weight $H_n$ only.
Sweeping the circle adds \emph{spatial resolution on $S^1$}: shell $n$ is an
$n$-way partition of the same compact carrier, not one more real height.
Divergent $H_n$ supplies depth; $2\pi k/n$ supplies angular resolution
(\leanid{harmonic\_depth\_diverges}, \texttt{HarmonicSweepCountingAxis}).
Each slot carries primitive-root phase on $S^1$
(\leanid{shell\_slot\_phase\_eq\_exp\_angle}) and rolls to
$\texttt{stripRollingMap}(2\pi k/n)$
(\leanid{strip\_rolling\_at\_slot\_eq\_lift}).
Head/tail reflection is antipodal, $\theta\mapsto\theta+\pi$
(\leanid{head\_tail\_reflect\_rolling\_antipodal}); reflections quotient slots,
they do not enumerate them.

\paragraph{$j$--$k$ unit-circle zero readout.}
On the line the Hopf lift is a point $(j,k)=(\cos t,\sin t)$ on $S^1$
(\leanid{hopfJKCirclePoint}).  The phase readout uses the \emph{coordinate
product} in place of a bare $\mathbb{C}$ generator:
$e^{\pi i j k}$ (\leanid{hopfJKUnitCirclePhase}); amplitude is the $45^\circ$
projection $j/\sqrt2+k/\sqrt2$ (\leanid{hopfJKCriticalAmplitude}).
\textbf{Unconditional:} twiddle vanishes iff $j+k=0$ on the circle
(\leanid{hopf\_jk\_twiddle\_vanishes\_iff\_cos\_sin}).
\textbf{Conditional:} $\zeta(\tfrac12+it)=0$ iff the amplitude vanishes under
\texttt{RollingZetaIdentificationAtCriticalLine}
(\leanid{zeta\_zero\_iff\_hopf\_jk\_twiddle}).
Harmonic-shell slots at angle $2\pi k/n$ inherit the same readout
(\leanid{harmonic\_shell\_balance\_iff\_jk\_product\_phase}).

\paragraph{Euler prime product counting on the circle.}
On $s=\tfrac12+it$ each factor splits as $n^{-s}=n^{-1/2}\cdot e^{-it\log n}$
(\leanid{critical\_line\_prime\_power\_polar}): modulus $n^{-1/2}$ carries no
angle; the unit phase $\texttt{linePhase}\,n\,t$ lives on the Hopf circle.
Multiplication in $\mathbb{N}$ is phase multiplication
(\leanid{euler\_phase\_multiplicative\_at\_slot}).  The natural counter is not a
height scan $t\in\mathbb{R}$ but a triple index
\[
  (\text{shell }n,\ \text{arc }k,\ \text{prime }p)
  \quad\longmapsto\quad
  \texttt{linePhase}\,p\,(2\pi k/n),
\]
together with the $j$--$k$ amplitude at that arc
(\leanfile{Hqiv/Story/S3HopfJKEulerPrimeCircleCounting.lean}).
Two distinct prime phases \emph{pin} the circle point
(\leanid{two\_prime\_phases\_pin\_circle\_angle}); one prime leaves only a
$2\pi/\log p$ ghost lattice.  Under the rolling bridge,
$\zeta=0$ at a slot angle iff the prime-log twiddle vanishes there
(\leanid{zeta\_zero\_iff\_shell\_prime\_twiddle}).

\textbf{Honesty.}  $\zeta(s)=0$ on the line as cover balance still requires
\texttt{RollingZetaIdentificationAtCriticalLine}
(\leanid{zeta\_zero\_iff\_cover\_balance}); the counting axis $(n,k)$ is
unconditional (\leanid{harmonicShellZeroCountingBundle}).

\subsection{45° readout on the whole strip}

The $45^\circ$ projection amplitude on the lift is
\[
  \texttt{criticalProj}(\texttt{stripAnalyticLift}(\sigma,t))
  = \frac{f(\sigma)+r(\sigma)(\cos t+\sin t)}{\sqrt2}
\]
(\leanid{strip\_analytic\_critical\_proj}).
Cancellation is the balanced condition
$f(\sigma)+r(\sigma)(\cos t+\sin t)=0$
(\leanid{strip\_analytic\_cancellation\_iff}).
The Fourier twiddle
$\texttt{stripAnalyticTwiddle}(\sigma,t)=e^{it}\cdot\texttt{criticalProj}(\cdots)$
extends the critical-line rolling twiddle to all $\sigma$
(\leanid{strip\_analytic\_twiddle\_on\_line}).

\subsection{Conditional zero classification}

\textbf{Honesty.} As in \S\ref{sec:real-line}, identifying $\zeta(s)$ with the
lifted residual requires \texttt{ZetaEqualsS3ResidualAt} and the sample
hypothesis \texttt{StripAnalyticLiftMatches} ($P.\texttt{coords}=\texttt{stripPointLift}(s)$).

\begin{theorem}[Whole-strip cancellation form]\label{thm:strip-cancel}
Under \texttt{StripAnalyticLiftMatches} and \texttt{ZetaEqualsS3ResidualAt},
\[
  \zeta(s)=0
  \;\Longleftrightarrow\;
  f(\sigma)+r(\sigma)(\cos t+\sin t)=0
\]
on the open strip (\leanid{zeta\_zero\_iff\_strip\_analytic\_cancellation}).
\end{theorem}

\begin{theorem}[Path C/E on the cylinder]\label{thm:strip-pathC}
Under the same hypotheses, for $0<\sigma<1$,
\[
  \zeta(s)=0
  \;\Longleftrightarrow\;
  \texttt{ZeroProducingOrbit}(\texttt{stripPointLift}(s))
  \;\wedge\;
  \texttt{PathEStripHolonomyCloses}(s)
\]
(\leanid{analytic\_strip\_pathC\_classification}).
On $\sigma=\tfrac12$ this recovers the rolling classification of
\leanfile{Hqiv/Story/S3StripRollingProjection.lean}
(\leanid{analytic\_strip\_rolling\_classification\_on\_line}).
\end{theorem}

The capstone \texttt{InteriorAssemblyNonzeroAtNontrivialZerosOffLine} is
unchanged: the lift supplies \emph{geometry} on $\sigma\times S^1$; the
identification bridge remains RH-hard.

\section{Regional closed form for $\zeta(s)$}\label{sec:closed-form}

Table~\ref{tab:regional} collects the proved regional formulas.  Every row is
machine-checked against Mathlib in \leanfile{Hqiv/Story/S3ZetaClosedForm.lean}
and \leanfile{Hqiv/Story/S3SO4ZetaProjectionClosedForm.lean}.

\begin{table}[h]
\centering
\caption{Regional closed forms for $\zeta(s)$ from the SO(4) projection (all rows proved).}
\label{tab:regional}
\begin{tabular}{@{}lll@{}}
\toprule
Region & Channel & Closed form \\
\midrule
$\Re s>1$ & even / Dirichlet & $\displaystyle\sum_{n\ge0}(n+1)^{-s}$ \\
$0<\Re s<1$ & odd / sin--cos--$\Gamma$ & FE assembly (\texttt{oddStripChannel}) \\
$s=2k$, $k\ge1$ & even $\pi$-sector & $\displaystyle(-1)^{k+1}\frac{2^{2k-1}\pi^{2k}B_{2k}}{(2k)!}$ \\
$s=-k$ & Bernoulli & $\displaystyle-\frac{B'_{k+1}}{k+1}$ \\
odd $s\ge3$ & odd continuation & no elementary closed form \\
\bottomrule
\end{tabular}
\end{table}

\subsection{Right half-plane}

\begin{theorem}[Shell / Dirichlet sum]
For $\Re s>1$,
\[
  \zeta(s)=\sum_{n=0}^{\infty}(n+1)^{-s},
\]
the HQIV shell ladder with shift $(n+1)$
(\leanid{riemannZeta\_dirichlet\_closed\_form}).
\end{theorem}

This is the even harmonic--Euler sector: the harmonic--prime--$\zeta$ path proves
the prime product and $\Lambda=-\zeta'/\zeta$ on the same region
(\leanfile{Hqiv/Story/S3HarmonicPrimeZetaPath.lean}).

\subsection{Open critical strip}

\begin{theorem}[Functional-equation closed form]\label{thm:fe-strip}
For $0<\Re s<1$,
\begin{equation}\label{eq:fe-strip}
  \zeta(s)
  =2\,(2\pi)^{-(1-s)}\,
  \Gamma(1-s)\,
  \cos\frac{\pi(1-s)}{2}\,
  \zeta(1-s).
\end{equation}
\end{theorem}

Equation~\eqref{eq:fe-strip} is Mathlib's \texttt{riemannZeta\_one\_sub}, repackaged
as \texttt{oddStripChannel} $=\zeta(s)$ on the strip
(\leanfile{Hqiv/Story/S3SO4InteriorWitness.lean}).

\subsection{Even positive integers}

\begin{theorem}[$\pi$-sector values]
For $k\ge1$,
\[
  \zeta(2k)=(-1)^{k+1}\frac{2^{2k-1}\pi^{2k}B_{2k}}{(2k)!}.
\]
In particular $\zeta(2)=\pi^2/6$ and $\zeta(4)=\pi^4/90$.
\end{theorem}

At these points $\mathcal{S}(2k)=0$: the sine/cosine rotation has completed full
$\pi$-turns and only the even $\pi$-sector residue remains---consistent with
Theorem~\ref{thm:even-cancel}.

\subsection{Negative integers}

\begin{theorem}[Bernoulli slot]
For $k\in\mathbb{N}$,
\[
  \zeta(-k)=-\frac{B'_{k+1}}{k+1}.
\]
\end{theorem}

\subsection{Master projection formula}

\begin{definition}
\texttt{zetaSO4Projected}$(s):=\zeta(s)$ with the regional characterisations of
Table~\ref{tab:regional} as proved corollaries.
\end{definition}

Off the critical line on the open strip,
\[
  \zeta(s)=\texttt{so4EquatorProject}(s)\cdot\texttt{interiorStripH}(s),
  \qquad
  \texttt{so4EquatorProject}(s)=\frac{2\sigma-1}{\sqrt2},
\]
when $0<\sigma<1$ and $\sigma\neq\tfrac12$
(\leanid{zetaSO4Projected\_eq\_sin\_cos\_channel\_off\_line}).

\section{Open-strip assembly and the critical line}\label{sec:real-line}

The open-strip functional-equation assembly~\eqref{eq:fe-strip} exposes three
features that look singular when read naively:
\begin{enumerate}[label=(\alph*),nosep]
\item $\Gamma(1-s)$ has poles at $s=1,2,3,\ldots$;
\item $\cos(\pi(1-s)/2)$ has zeros at odd integers;
\item negative-even \emph{trivial} zeros of $\zeta$ sit at $s=-2,-4,\ldots$
\end{enumerate}

The machine-checked classification treats these as \textbf{regional bookkeeping},
not as contradictions:
\begin{itemize}[nosep]
\item \textbf{Negative-even trivial zeros} are the only negative-integer zeros;
      nontrivial zeros cannot occupy negative integer slots
      (\leanfile{Hqiv/Story/S3FESlotDischarge.lean},
      \leanid{nontrivial\_zero\_fe\_slot}).
\item \textbf{Open strip} points avoid $\Re s>1$ and $\Re s<0$ by Mathlib
      nonvanishing and FE reflection
      (\leanfile{Hqiv/Story/S3DeltaOrbitOffStrip.lean}).
\item \textbf{$\Gamma$ poles} on the strip are paired with cos zeros in the
      standard FE; they mark the even-sector continuation, not a complex off-line
      zero locus.
\end{itemize}

\section{Orbit cancellation vs.\ pointwise vanishing}\label{sec:orbit-gap}

\leanfile{Hqiv/Story/S3OrbitVsPointwiseGap.lean} separates three layers:
\begin{enumerate}[label=(\roman*),nosep]
\item \textbf{Orbit cancellation.}  For every $\sigma$,
      $\texttt{rot45Free}(\sigma)+\texttt{rot45Free}(1-\sigma)=0$
      (\leanid{orbit\_free\_sum\_cancels}); this holds even off
      $\Re s=\tfrac12$ and cannot exclude off-line $\zeta$ zeros by itself.
\item \textbf{Pointwise vanishing.}  A single free coordinate vanishes iff
      $\sigma=\tfrac12$ (\leanid{pointwise\_free\_zero\_iff\_on\_line}).
\item \textbf{RH equivalence.}  \texttt{AllNontrivialZerosOnLine}
      $\Leftrightarrow$ \texttt{RiemannHypothesis}
      (\leanid{allNontrivialZerosOnLine\_iff\_RiemannHypothesis}).
\end{enumerate}

\subsection{FE + Schwarz quadruplet}\label{sec:quadruplet}

Nontrivial zeros satisfy the FE pair $\zeta(1-s)=0$ unconditionally
(\leanfile{Hqiv/Story/S3ZeroQuadrupletOrbit.lean},
\leanid{zeta\_zero\_fe\_pair}).  The full Schwarz quadruplet
\[
  \{s,\,1-s,\,\overline{s},\,1-\overline{s}\}
\]
is packaged as \texttt{zetaZeroQuadruplet}; all four points carry zeros when
the conjugation identity $\zeta(\overline{t})=\overline{\zeta(t)}$ is available
(\leanid{zeta\_zero\_quadruplet}).  On $\Re s=\tfrac12$ the reflections
$\overline{s}$ and $1-s$ coincide, collapsing the quadruplet to an FE pair---the
same collapse as orbit vs.\ pointwise on the $\sigma$ channel.

\begin{remark}[Geometric reading]
The ``infinity'' from $\Gamma$ and the ``zero'' from negative-even slots are
two faces of the same projection geometry: they clear the plane of spurious
complex-axis candidates and leave the \emph{pointwise} equator condition
$\texttt{rot45Free}=0\Leftrightarrow\sigma=\tfrac12$ as the only balanced
$\sigma$-only dependence.  RH asks that actual $\zeta$ zeros satisfy this stronger
pointwise condition, not merely orbit cancellation along $\sigma\leftrightarrow
1-\sigma$.
\end{remark}

\subsection{Tangent law of the orbit: factorization angles and the unimodular
locus}\label{sec:tangent-orbit}

The strip closed form carries a \emph{factorization angle} $\pi s/2$: the sine/cosine
rotation slots are $\sin(\pi s/2)$ and $\cos(\pi s/2)$ and their ratio
$T(s)=\tan(\pi s/2)$ is the projection formula
(\texttt{zetaAxisRotationRatio}).
\leanfile{Hqiv/Story/S3TangentOrbitCriticalLine.lean} proves the orbit law
for that angle under the FE + Schwarz quadruplet --- all of it
unconditional projection geometry, no RH input.

\paragraph{The orbit acts on the tangent by inversion and conjugation.}
FE reflection swaps the slots exactly
($\sin\!\big(\pi(1-s)/2\big)=\cos(\pi s/2)$,
\leanid{proj\_sin\_one\_sub}), so
\[
  T(1-s) = T(s)^{-1},
  \qquad
  T(\overline{s}) = \overline{T(s)},
  \qquad
  T(1-\overline{s}) = \overline{T(s)}^{-1}
\]
(\leanid{projTangent\_one\_sub}, \leanid{projTangent\_schwarz},
\leanid{projTangent\_quadruplet}), with the orbit reciprocity
$T(1-s)\,T(s)=1$ on the strip
(\leanid{projTangent\_orbit\_reciprocity}).  The quadruplet
$\{s,1-s,\overline{s},1-\overline{s}\}$ reads out through the tangent as
$\{T,\,T^{-1},\,\overline{T},\,\overline{T}^{-1}\}$: the group
$\{1,\mathrm{inv},\mathrm{conj},\mathrm{inv}\circ\mathrm{conj}\}$.

\paragraph{The critical line is exactly the unimodular locus.}
The engine is the exact modulus identity
\[
  \big\lVert\sin(\pi s/2)\big\rVert^2-\big\lVert\cos(\pi s/2)\big\rVert^2
  \;=\; -\cos\!\big(\pi\,\Re s\big)
\]
(\leanid{normSq\_proj\_sin\_sub\_cos}): the $\Im s$ content rides in common
$\cosh/\sinh$ factors that cancel in the difference --- the $\sigma$-blindness
of \S\ref{sec:sigma-scope}, reappearing at the tangent level.  Consequently,
on the open strip,
\[
  \lVert T(s)\rVert^2 = 1 \;\Longleftrightarrow\; \Re s=\tfrac12,
  \qquad
  \lVert T(s)\rVert^2 < 1 \;\Longleftrightarrow\; \Re s<\tfrac12,
  \qquad
  \lVert T(s)\rVert^2 > 1 \;\Longleftrightarrow\; \Re s>\tfrac12
\]
(\leanid{normSq\_projTangent\_eq\_one\_iff} and companions): the tangent
modulus is a \emph{complete invariant} for the critical line and sorts the
strip around it.

\paragraph{Orbit collapse $\Leftrightarrow$ critical line.}
Inversion equals conjugation on the tangent --- $T^{-1}=\overline{T}$,
collapsing the quadruplet tangent orbit from four values to two --- precisely
on the critical line (\leanid{projTangent\_inv\_eq\_conj\_iff}).  And the tie
to the $\zeta$ factorization is exact: the equator factor of
$\zeta=h\cdot(2\sigma-1)/\sqrt2$ vanishes exactly where the tangent is
unimodular,
\[
  \texttt{so4CriticalFactor}(s)=0
  \quad\Longleftrightarrow\quad
  \lVert\tan(\pi s/2)\rVert^2=1
\]
(\leanid{so4CriticalFactor\_zero\_iff\_unimodular\_tangent}).  The critical
line is thus characterized three ways at once --- vanishing equator factor,
unimodular factorization tangent, and coincidence of the two orbit
symmetries --- all proved to be the same locus.

\subsection{Zeros as spectral cancellations of the Euler product}
\label{sec:spectral-cancellation}

The reading that drives this entire construction --- \emph{the zeros are
exactly the spectral cancellations that give the Euler prime product} ---
has three provable faces, and
\leanfile{Hqiv/Story/S3EulerSpectralCancellation.lean} proves all three.
None of them requires RH.

\paragraph{In the product region the prime phases provably cannot cancel.}
For $\Re s>1$ every Euler factor is individually nonvanishing,
$\lVert p^{-s}\rVert=p^{-\Re s}<1$ so $(1-p^{-s})^{-1}\neq0$
(\leanid{euler\_factor\_ne\_zero}); the product of these phases converges to
$\zeta(s)$ \emph{and} the value is nonzero
(\leanid{euler\_product\_region\_no\_cancellation}, packaging Mathlib's
\texttt{riemannZeta\_eulerProduct\_hasProd} with the nonvanishing).  Where
the spectral product is a genuine convergent product, cancellation can never
\emph{complete}: a zero would contradict the product representation itself.

\paragraph{Zeros live exactly on the breakdown locus.}
Contrapositively, $\zeta(s)=0$ forces $\Re s<1$ --- impossible anywhere the
product region extends, including the closed boundary $\Re s=1$
(\leanid{zero\_forces\_product\_breakdown}).  Combined with the strip
confinement of \S\ref{sec:real-line}, every nontrivial zero sits in the open
strip $0<\Re s<1$ (\leanid{spectral\_cancellation\_in\_breakdown\_locus}) ---
precisely the region where the prime spectral series no longer converges
absolutely.  The zero set lives on the breakdown locus of the Euler product
and nowhere else.

\paragraph{Each zero is a resonance pole of the inverse spectral response.}
On the product region the prime spectrum is \emph{tame}: the von Mangoldt
series converges and equals the log-derivative,
$\sum_n\Lambda(n)n^{-s}=-\zeta'(s)/\zeta(s)$
(\leanid{spectral\_series\_tame}).  The \emph{inverse} spectral response
$1/\zeta$ --- the un-inverted prime-phase product $\prod_p(1-p^{-s})$
continued to the strip --- blows up at every zero:
\[
  \big\lVert \zeta(s)^{-1}\big\rVert \;\longrightarrow\; \infty
  \qquad\text{as } s\to\rho \text{ along the nonvanishing locus,}
\]
unconditionally for any zero $\rho\neq1$
(\leanid{inverse\_spectral\_response\_blows\_up\_at\_zero}); the
punctured-neighbourhood version carries the classical isolated-zero fact as
an explicit hypothesis rather than smuggling it in
(\leanid{inverse\_spectral\_response\_blows\_up\_punctured}).  A completed
cancellation is a \emph{resonance}: the spectral measure of the primes
diverges exactly there.

\paragraph{Capstone: the spectral-cancellation profile.}
\leanid{spectral\_cancellation\_profile} packages, for any nontrivial zero
$\rho$: (i) open-strip confinement (breakdown region); (ii) impossibility in
the product region; (iii) the resonance blow-up of $\lVert1/\zeta\rVert$ at
$\rho$; and (iv) the quadruplet tangent orbit
$\{T,T^{-1},\overline{T},\overline{T}^{-1}\}$ of
\S\ref{sec:tangent-orbit} attached to the factorization angle $\pi\rho/2$.
A nontrivial zero \emph{is} a completed spectral cancellation carrying the
full orbit data.  What RH adds --- and only RH --- is that every such
cancellation sits on the unimodular-tangent locus $\lVert T\rVert=1$.

\subsection{Resonance channeling and the SO(4) spectral-line transformer}
\label{sec:resonance-channeling}

\leanfile{Hqiv/Story/S3SpectralResonanceChanneling.lean} combines the three
proved layers --- tangency, equator factor, and spectral cancellation --- into the
sharpest forcing statements available without assuming RH.

\paragraph{The cancellation is collective.}
A single prime line never cancels on the entire right half-plane: for
$\Re s>0$ (in particular on the whole strip), $1-p^{-s}\neq0$
(\leanid{euler\_factor\_base\_ne\_zero\_right\_half}), so every finite
spectral truncation $\prod_{p\in F}(1-p^{-s})$ is nonzero
(\leanid{finite\_spectral\_truncation\_ne\_zero}).  A zeta zero needs all
primes at once; no finite interference produces it.

\paragraph{The two-channel dichotomy.}
On the strip $\zeta=h\cdot(2\sigma-1)/\sqrt2$ off the line, so at any
nontrivial zero the collective cancellation must be absorbed by exactly one
of two channels (\leanid{resonance\_channeling\_dichotomy}):
\[
  \underbrace{\Re\rho=\tfrac12:\;
    \texttt{so4CriticalFactor}(\rho)=0,\;\lVert T(\rho)\rVert^2=1}_{
    \text{geometric equator channel absorbs}}
  \qquad\text{or}\qquad
  \underbrace{\Re\rho\neq\tfrac12:\;
    \texttt{interiorStripH}(\rho)=0}_{
    \text{free assembly forced to vanish}}
\]
(\leanid{offline\_zero\_forces\_assembly\_vanish}).  RH is exactly the
statement that the second channel never fires, and this is now a proved
\emph{equivalence}, not merely an implication:
\leanid{geometric\_channel\_absorption\_iff\_RH} gives
$\texttt{InteriorAssemblyNonzeroAtNontrivialZerosOffLine}\;
\texttt{interiorStripH}\Leftrightarrow\text{RH}$.

\paragraph{The SO(4) spectral-line transformer.}
For a given $n$, the transformer is the spectral line
$\texttt{so4SpectralLine}\,n\,s=n^{-s}$, with proved laws: the modulus law
$\lVert n^{-s}\rVert=n^{-\Re s}$ ($\sigma$-blindness at the spectrum level,
\leanid{so4SpectralLine\_norm}); the unconditional \emph{FE pairing}
\[
  \lVert n^{-s}\rVert\cdot\lVert n^{-(1-s)}\rVert=\frac1n
  \qquad\text{for \emph{all} } s
\]
(\leanid{so4SpectralLine\_fe\_product}); the \emph{geometric-mean condition}
--- the two paired lines are equal iff $\Re s=\tfrac12$
(\leanid{so4SpectralLine\_geometric\_mean\_iff}); the
\emph{square-root-weight condition} $\lVert n^{-s}\rVert^2=1/n
\Leftrightarrow\Re s=\tfrac12$ (\leanid{so4SpectralLine\_sq\_weight});
multiplicativity $(ab)^{-s}=a^{-s}b^{-s}$
(\leanid{so4SpectralLine\_mul}); and the prime generation of every line ---
every $n\ge2$ is a multiple of a prime and its line factors through that
prime's line (\leanid{spectral\_line\_prime\_factorization}).

\paragraph{Main conditional hypothesis.}
\leanid{RH\_iff\_zero\_spectral\_weights}: \textbf{RH is equivalent to the
statement that at every nontrivial zero, every spectral line carries exactly
square-root weight} $\lVert n^{-\rho}\rVert^2=1/n$.  ``For every prime, the
zero sits at its square-root pole weight'' --- made exact and proved
RH-equivalent.  And \leanid{zero\_line\_characterizations} shows the three
equivalent conditions agree at every zero: vanishing equator factor $\Leftrightarrow$
unimodular tangent $\Leftrightarrow$ all lines at square-root weight
$\Leftrightarrow\Re\rho=\tfrac12$.  The unconditional content is a genuine
\emph{reduction}: RH is now a single-channel statement --- the free assembly
$h$ never absorbs a collective cancellation --- with every geometric equivalent condition
proved to name the same line.

\subsection{Vertical axis rigidity: the orbits cannot wobble}
\label{sec:axis-rigidity}

\leanfile{Hqiv/Story/S3VerticalAxisRigidity.lean} formalizes the rigidity
intuition: any ``curvature'' of the zero locus --- a second symmetry axis, a
tilted or bent critical line --- would force zeros \emph{out of the strip,
off to infinity}.

\paragraph{Two mirrors make a translation.}
The zero set is unconditionally closed under the FE point-reflection
$\rho\mapsto1-\rho$ (\leanid{nontrivial\_zero\_fe\_closed} --- FE only, no
Schwarz hypothesis needed).  Suppose it were \emph{also} closed under the
mirror $\texttt{verticalReflect}\,\sigma_0:s\mapsto2\sigma_0-\overline{s}$
about some axis $\Re s=\sigma_0$.  Composing the two symmetries twice yields
the translation
\[
  s\;\longmapsto\;s+(2-4\sigma_0),
\]
nontrivial unless $\sigma_0=\tfrac12$; iterating it marches any zero
monotonically in $\Re$, straight out of the open strip --- contradicting the
proved strip confinement.  Hence
(\leanid{strip\_set\_unique\_vertical\_axis}, abstract;
\leanid{zeta\_zero\_unique\_vertical\_axis}, instantiated): \textbf{the
nontrivial zero set admits at most one vertical mirror axis, and it must be
$\Re s=\tfrac12$.}  The only hypothesis carried by the zeta instantiation is
nonemptiness of the zero set (Hardy 1914, classical, not yet in Mathlib).

\paragraph{Supporting rigidity facts (all unconditional).}
The fixed-point set of any vertical mirror is exactly its axis --- a
straight line, never a curve (\leanid{verticalReflect\_fixed\_iff}); the
mirror about $\tfrac12$ is the FE+Schwarz composite $1-\overline{s}$
(\leanid{verticalReflect\_half\_eq\_one\_sub\_schwarz}).  If a point and its
FE partner $1-s$ share a vertical line, that line is $\Re s=\tfrac12$
(\leanid{fe\_pair\_same\_vertical\_line\_iff}): any two symmetry-partner
points on one line force \emph{the} line.  And the candidate locus itself
has no curvature anywhere: as a subset of all of $\mathbb{C}$, the
square-root-weight locus is exactly the straight vertical line
$\{\Re s=\tfrac12\}$ (\leanid{sq\_weight\_locus\_is\_vertical\_line}), swept
by the vertical flow $s\mapsto s+it$ with no transverse drift
(\leanid{sq\_weight\_locus\_vertical\_flow}).

\paragraph{Honest scope.}
``Two points on the line force all points onto the line'' is not a valid
inference classically.  What is proved: (i) two symmetry-partner points on a
common vertical line pin that line to $\tfrac12$; (ii) the candidate locus
is a perfect straight line, not a curve; (iii) the zero set cannot support a
second mirror axis without leaking out of the strip.  Promoting ``the unique
possible axis is $\tfrac12$'' to ``every zero lies on it'' is exactly RH.

\subsection{Same-height orbit collapse: ``one zero per height'' is exactly
RH}\label{sec:same-height}

\leanfile{Hqiv/Story/S3SameHeightOrbitCollapse.lean} scopes the two-part
equivalence --- \emph{collectivity pins the axis; two zeros at the same
height are off the board} --- and extracts the genuine theorem inside it.

\paragraph{Collectivity does not pin the axis.}
The cancellation at every zero is provably collective, but at \emph{any}
point of the strip the all-primes weight family is the perfectly consistent
$\lVert n^{-s}\rVert^2=n^{-2\sigma}$
(\leanid{spectral\_weights\_consistent\_at\_any\_sigma}): nothing diverges
or contradicts off the line.  What distinguishes $\sigma=\tfrac12$ is solely
the square-root normalization $n^{-1}$ --- and that normalization \emph{is}
RH (\leanid{RH\_iff\_zero\_spectral\_weights}).

\paragraph{The same-height partner and the merge law.}
The quadruplet member $1-\overline{s}$ sits at the \emph{same height},
$\Im(1-\overline{s})=\Im s$ (\leanid{one\_sub\_schwarz\_same\_height}), and
merges with $s$ exactly on the critical line
(\leanid{same\_height\_partner\_merges\_iff} --- the mirror fixed-point law
again).  An off-line zero would therefore carry a \emph{distinct} second
zero at its own height; an on-line zero carries none.

\paragraph{``One zero per height'' $\Leftrightarrow$ RH.}
Declaring the straddling case off the board is not a proof of RH --- it is an
exact reformulation of RH, proved as an equivalence in two strengths:
\begin{itemize}[nosep]
\item \leanid{RH\_iff\_height\_pins\_re} (\textbf{FE only, unconditional}):
      RH $\Leftrightarrow$ any two nontrivial zeros with the same
      $\lvert\Im\rvert$ have the same $\Re$.  The backward direction needs
      no Schwarz input: the FE partner $1-\rho$ already supplies a zero of
      equal $\lvert\Im\rvert$ with reflected real part, forcing
      $\Re\rho=1-\Re\rho$.
\item \leanid{RH\_iff\_unique\_zero\_per\_height} (Schwarz identity carried
      as the explicit hypothesis, in the style of
      \leanid{zeta\_zero\_quadruplet}): RH $\Leftrightarrow$ each height
      carries at most one nontrivial zero.
\end{itemize}
The orbit picture is correct: a quadruplet either straddles the line with
two distinct same-height members or collapses onto it.  Both directions are
machine-checked with exact equivalence.

\subsection{Polar projection collapse: the poles must land on one point}
\label{sec:polar-collapse}

\leanfile{Hqiv/Story/S3PolarProjectionCollapse.lean} formalizes the polar
reading: \emph{the poles of $S^3$ must project to a single point on the
complex plane for the zero to sit on the line} --- and resolves the worry
that ``SO(4) is only applicable if RH is true.''

\paragraph{The polar separation is the equator factor.}
At the height of $s$ the mirror presents two pole images, $s$ and
$\texttt{polarPartner}\,s=1-\overline{s}$, always at the same height
(\leanid{polarPartner\_same\_height}).  Their separation is horizontal,
real, and \emph{exactly} the equator factor:
\[
  \texttt{polarPartner}\,s-s \;=\; 1-2\sigma
  \;=\; -\sqrt2\cdot\texttt{so4CriticalFactor}(s)
\]
(\leanid{polar\_separation\_eq},
\leanid{polar\_separation\_eq\_twiddle}), with norm equal to twice the
distance to the critical line
(\leanid{polar\_separation\_eq\_twice\_deviation}).  The $45^\circ$ equator factor
factor is, verbatim, the displacement between the two pole projections: off
the line the SO(4) geometry does not break --- it \emph{reports} the polar
separation through the surviving equator channel, which is precisely the
off-line branch of the channeling dichotomy of
\S\ref{sec:resonance-channeling}.

\paragraph{Single-point projection $\Leftrightarrow$ RH.}
The two pole images coincide iff $\Re s=\tfrac12$
(\leanid{polar\_collapse\_iff\_on\_line}).  Single-point polar projection on
the zero set (\leanid{PolarProjectionCollapsesOnZeros}) is therefore not a
hidden presupposition of the SO(4) machinery but a well-formed proposition,
and it is RH --- proved as a exact-equivalence equivalence
(\leanid{polar\_collapse\_iff\_RH}), chaining to ``one zero per height''
given the Schwarz identity
(\leanid{polar\_collapse\_iff\_unique\_height}).  Single-point polar
projection at every zero is thus the fourth face of RH in this stack, equal
in strength to unimodular tangent, square-root spectral weights, and
one-zero-per-height.

\subsection{Zero--holonomy--Goldbach chain: every zero activates every
pair}\label{sec:zero-holonomy-chain}

\leanfile{Hqiv/Story/S3ZeroHolonomyGoldbachChain.lean} chains the
collectivity of the cancellation through the Goldbach parity bridge: since
all zeros are comprised of all primes, every zero carries the holonomy of
every Goldbach pair.

\paragraph{Activation and the AM--GM weight cap (unconditional).}
At any nontrivial zero $\rho$ and for any midpoint pair $p+q=2N$, both prime
lines are nonvanishing participants in the collective cancellation
(\leanid{zero\_activates\_pair}), and the joint line factors
multiplicatively through the pair.  AM--GM gives $pq\le N^2$
(\leanid{midpoint\_pair\_product\_le}), so on the critical line the joint
pair weight obeys
\[
  \lVert (pq)^{-s}\rVert^2=\frac1{pq}\;\ge\;\frac1{N^2}
\]
(\leanid{midpoint\_pair\_joint\_weight\_bound}): the \emph{additive}
half-slope datum $N$ caps the \emph{multiplicative} spectral weight of its
pairs.

\paragraph{Every zero contains the holonomy of every pair.}
The main chain theorem \leanid{zero\_contains\_pair\_holonomy} states, at any
nontrivial zero and for any midpoint pair: activation of both prime lines,
multiplicative factorization of the joint line, SO(4) midpoint slope
$\tfrac12$, and the Hopf-fiber holonomy support certificate (integrable
shell with \texttt{SO8AdmissibleHolonomy} $\Delta$/G$_2$/triality fields)
--- all unconditional.  Under \texttt{GoldbachParity}, every midpoint
$N\ge2$ admits a holonomy-supported pair
(\leanid{parity\_gives\_holonomy\_support}).

\paragraph{The bridge in geometric-channel language.}
The SO(8) half-slope bridge now rewrites with both channels geometric:
\begin{align*}
  \texttt{SO8ProjectedHalfSlopeBridge}\,2
  &\;\Longleftrightarrow\;
  \text{(square-root spectral weights at every zero)}\wedge\text{parity}\\
  &\;\Longleftrightarrow\;
  \text{(single-point polar projection at every zero)}\wedge\text{parity}
\end{align*}
(\leanid{bridge\_iff\_spectral\_weights\_and\_parity},
\leanid{bridge\_iff\_polar\_collapse\_and\_parity}).  The chain does not
prove either side --- the bridge remains exactly RH $\wedge$ Goldbach
with exact equivalence --- but both sides now speak the same SO(4)/holonomy
dialect: zeta-side polar collapse, prime-side holonomy-supported midpoint
pairs.

\subsection{TUFT nested frame tower: odd spheres and the harmonic radius
ladder}\label{sec:tuft-tower}

\leanfile{Hqiv/Story/S3TuftNestedFrameTower.lean} formalizes what the TUFT
nested fibration --- odd spheres $S^{2n+1}$ for every $n$ --- buys.  At
level $N$ the first $N$ spectral lines assemble into the frame vector
$(1^{-s},\dots,N^{-s})\in\mathbb{C}^N$, whose normalized version lives on
$S^{2N-1}$; the structure is carried by the squared radius
$\texttt{spectralFrameNormSq}\,N\,s=\sum_{n\le N}\lVert n^{-s}\rVert^2$.

\paragraph{Tower structure (unconditional).}
The levels nest with the new line's weight as increment
(\leanid{spectralFrameNormSq\_succ}); every level is $\sigma$-blind, so the
vertical flow moves the frame on a sphere of \emph{fixed} radius
(\leanid{spectralFrameNormSq\_eq\_of\_re\_eq}); the radius never
degenerates (\leanid{spectralFrameNormSq\_pos}).

\paragraph{The harmonic radius law.}
The rigidity payoff is a \emph{ladder of equivalent conditions}: for every level
$N\ge2$,
\[
  \texttt{spectralFrameNormSq}\,N\,s = H_N
  \quad\Longleftrightarrow\quad
  \Re s=\tfrac12
\]
(\leanid{frame\_radius\_eq\_harmonic\_iff}), where $H_N$ is the harmonic
backbone of the closure story.  The critical line is the unique vertical
line on which the nested sphere radii reproduce the harmonic ladder ---
\emph{the divergent series that forced $\Delta$ and closed
$\mathfrak{so}(3)+\Delta$ to $\mathfrak{so}(4)$ in the first place}.  The
loop closes: on the line the tower rebuilds the harmonic divergence
(\leanid{frame\_tower\_diverges\_on\_line}, via Mathlib's divergence of the
harmonic series), and the closure inequality $H_N\le K(N)$ reads as a
curvature bound on the nested sphere radii
(\leanid{frame\_radius\_le\_curvature\_on\_line}).

\paragraph{What it buys --- in classical terms.}
\leanid{RH\_iff\_zero\_frame\_tower\_harmonic}: RH $\Leftrightarrow$ at
every nontrivial zero the \emph{whole tower} carries harmonic radii --- an
infinite ladder of equivalent reformulations, one per level, all collapsing
to the same bit.  The tower adds consistency rigidity (infinitely many
independent conditions provably naming one line, with the closure backbone as
reference values), not a new forcing mechanism: each level is
  RH-equivalent, not RH-implying.

\subsection{The Diophantine transformer: the operator behind the
tower}\label{sec:diophantine-transformer}

The tower, the spectral lines, and the FE pairing are all shadows of one
concrete object.  \leanfile{Hqiv/Story/S3DiophantineTransformer.lean}
defines the level-$N$ \emph{Diophantine transformer}
\[
  D_N(s) \;=\; \mathrm{diag}\bigl(1^{-s},\,2^{-s},\,\dots,\,N^{-s}\bigr)
  \;\in\; M_N(\mathbb{C}),
\]
an integer-indexed, prime-generated diagonal operator on $\mathbb{C}^N$
(entries are multiplicative in the index, so the primes generate
everything), and proves its laws unconditionally.

\paragraph{Trace ladder $\to\zeta$.}
$\operatorname{tr}D_N(s)$ is the $N$-th Dirichlet partial sum
(\leanid{diophantineTransformer\_trace}), and on the product region
$\Re s>1$ the trace ladder converges:
$\operatorname{tr}D_N(s)\to\zeta(s)$
(\leanid{diophantineTransformer\_trace\_tendsto}).  The zeta function
\emph{is} the trace of the transformer tower.

\paragraph{HS norm $=$ tower radius.}
The squared Hilbert--Schmidt norm of $D_N(s)$ is exactly
$\texttt{spectralFrameNormSq}\,N\,s$
(\leanid{diophantineTransformer\_hs\_normSq}): the TUFT frame radius of
\S\ref{sec:tuft-tower} is the operator's norm ladder, and the harmonic
radius law becomes ``$D_N$ is harmonically normalized
($\lVert D_N\rVert_{\mathrm{HS}}^2=H_N$) exactly on the critical line.''

\paragraph{FE composition is a fixed rational operator.}
For \emph{every} $s$,
\[
  D_N(s)\cdot D_N(1-s)
  \;=\;
  \mathrm{diag}\bigl(1,\tfrac12,\tfrac13,\dots,\tfrac1N\bigr)
\]
(\leanid{diophantineTransformer\_fe\_product}) --- the $s$-dependence
cancels completely and what remains is a pure Diophantine object with
rational entries, whose trace is the harmonic backbone $H_N$
(\leanid{feProduct\_trace\_eq\_harmonic}).  The functional equation, seen
through the transformer, is the statement that the two reflected operators
always compose to the \emph{integer harmonic operator}.

\paragraph{The adjoint law names the line.}
The genuinely new rigidity is at the level of adjoints: for any $n\ge2$,
\[
  n^{-(1-s)}=\overline{n^{-s}}
  \quad\Longleftrightarrow\quad
  \Re s=\tfrac12
\]
(\leanid{transformer\_entry\_adjoint\_iff}), hence for $N\ge2$
\[
  D_N(1-s)=D_N(s)^{\dagger}
  \quad\Longleftrightarrow\quad
  \Re s=\tfrac12
\]
(\leanid{diophantineTransformer\_adjoint\_iff}).  The FE reflection agrees
with the operator adjoint precisely on the critical line, and there the
transformer is a \emph{scaled isometry} whose unitarity defect is again the
fixed harmonic diagonal:
$D_N(s)D_N(s)^{\dagger}=\mathrm{diag}(1/n)$
(\leanid{transformer\_scaled\_isometry\_on\_line}).

\paragraph{RH in operator language.}
\leanid{RH\_iff\_transformer\_adjoint\_at\_zeros}: RH $\Leftrightarrow$ at
every nontrivial zero and every level $N\ge2$,
$D_N(1-\rho)=D_N(\rho)^{\dagger}$.  This is the machine-checked
form of the Hilbert--P\'olya intuition the construction has been circling:
not ``there exists a self-adjoint operator whose spectrum is the zeros,''
but ``RH $\Leftrightarrow$ the concrete transformer family is
self-adjoint-across-FE at the zeros.''  As with every equivalent condition in this
paper, the law is an exact reformulation of RH, not a proof of RH.

\subsection{The octonionic associator channel: breaking the abelian
obstruction}\label{sec:associator-channel}

Diagonalization itself is the obstruction.
\leanfile{Hqiv/Story/S3OctonionicAssociatorChannel.lean} first proves the
saturation formally: the transformer family is \emph{abelian} and
\emph{normal},
\[
  D_N(s)\,D_N(s') = D_N(s')\,D_N(s),
  \qquad
  D_N(s)\,D_N(s)^{\dagger} = D_N(s)^{\dagger}\,D_N(s)
\]
(\leanid{diophantineTransformer\_mul\_comm},
\leanid{diophantineTransformer\_normal}).  A commuting normal family is
simultaneously diagonalizable; every invariant it carries factors through
the entries, and the entries are exactly the $\sigma$-blind /
RH-equivalent readouts catalogued above.  This is the machine-checked
``good reason'' the diagonal attack cannot go further.

\paragraph{The escape hatch is non-associativity.}
The HQIV stack carries one algebra that no commuting family can simulate:
the octonionic light-cone carrier, whose associator
$[x,y,z]=(xy)z-x(yz)$ is the same phase channel that supplies CP-odd
structure and Lorentz bookkeeping elsewhere in the framework.  On the
frozen Fano tables the witness is concrete:
\[
  [e_1,e_2,e_4] \;=\; (e_1e_2)e_4 - e_1(e_2e_4)
  \;=\; e_7e_4 - e_1(-e_5)
  \;=\; -e_3 - e_4 \;\neq\; 0,
\]
with squared norm $2$
(\leanid{octonionAssociator\_e1\_e2\_e4},
\leanid{octonionAssociatorNormSq\_e1\_e2\_e4}).

\paragraph{Spectral weights ride the torsion.}
The associator is trilinear in scalar weights
(\leanid{octonionAssociator\_smul}), so the prime spectral moduli couple
cleanly: weighting $(e_1,e_2,e_4)$ by
$\lVert a^{-s}\rVert,\lVert b^{-s}\rVert,\lVert c^{-s}\rVert$ defines the
\emph{associator channel} with closed form
\[
  \texttt{octAssociatorChannel}\;a\,b\,c\,s
  \;=\; 2\,(abc)^{-2\sigma}
\]
(\leanid{octAssociatorChannel\_eq}).  Three unconditional laws follow:
\begin{itemize}[nosep]
\item \textbf{Zeros cannot cancel the torsion}
      (\leanid{zero\_keeps\_associator\_torsion}): the channel is strictly
      positive at every point — in particular at every nontrivial zero.
      The zero annihilates the abelian assembly $\sum n^{-s}$; the
      non-associative channel survives untouched, carrying every prime's
      weight through the cancellation.
\item \textbf{Asymmetry condition}
      (\leanid{octAssociatorChannel\_eq\_iff}): the channel carries the
      square-root weight $2/(abc)$ exactly on $\Re s=\tfrac12$.
\item \textbf{FE pairing} (\leanid{octAssociatorChannel\_fe\_product}):
      channel$(s)\cdot$channel$(1-s)=4/(abc)^2$ for every $s$ — the
      associator analogue of $D_N(s)D_N(1-s)=\mathrm{diag}(1/n)$.
\end{itemize}
The Goldbach hook also extends: for a midpoint pair $p+q=2N$, the triple
$(p,q,2N)$ keeps the channel above $1/N^3$ on the line
(\leanid{midpoint\_triple\_channel\_floor}) — the AM--GM cap pushed into
the non-associative channel.

\paragraph{What the asymmetry buys --- in classical terms.}
\leanid{RH\_iff\_zero\_associator\_channel}: RH $\Leftrightarrow$ at every
nontrivial zero the channel carries exact square-root weight for every
triple.  The channel structurally breaks the abelian obstruction — it is
an orientation-sensitive three-slot invariant that no commuting diagonal
family can produce, and unlike the assembly it is provably \emph{active}
at every zero.  But its modulus coupling is still driven by $\Re s$
alone, so the condition collapses to the same single bit: another exact
reformulation.  What is genuinely new is the \emph{shape} of the
invariant (a non-associative three-form versus abelian diagonals) and the
unconditional positivity at zeros; converting the torsion's phase
content — rather than its modulus — into a forcing mechanism is the open
edge this section names.

\subsection{S\texorpdfstring{$^7$}{7} torsion cancellation: the SO(8)
twist}\label{sec:s7-torsion}

\leanfile{Hqiv/Story/S3OctonionS7TorsionCancellation.lean} keeps the full
\emph{torsion vector}
$\texttt{octTorsionVector}\,a\,b\,c\,s\in\mathbb{R}^8$ in the carrier and
asks where it cancels between FE partners.  The carrier's unit sphere is
$S^7$; the rotation group is SO(8), whose complete invariant on
$\mathbb{R}^8$ is the norm.

\paragraph{The S$^7$ shadow never moves.}
For every spectral parameter,
\[
  \texttt{octTorsionVector}\,a\,b\,c\,s
  \;=\;
  (abc)^{-\sigma}\,\cdot\,(-e_3-e_4)
\]
(\leanid{octTorsionVector\_eq}): a \emph{positive} radial weight times one
fixed carrier direction (\leanid{octTorsion\_single\_point\_on\_s7}).  The
entire spectral flow is radial; $S^7$ sees a single point, and all the
information sits in the weight $(abc)^{-\sigma}$.

\paragraph{Cancellation, chirality, orbits.}
The FE defect
$\texttt{octTorsionDefect}=\texttt{torsion}(s)-\texttt{torsion}(1-s)$ is
the weight gap $(abc)^{-\sigma}-(abc)^{-(1-\sigma)}$ times the fixed
direction (\leanid{octTorsionDefect\_eq}), is odd under $s\mapsto1-s$
(\leanid{octTorsionDefect\_antisymm}), and:
\begin{itemize}[nosep]
\item \textbf{cancels exactly on the line}
      (\leanid{octTorsionDefect\_zero\_iff}): torsion$(s)=$
      torsion$(1-s)$ $\Leftrightarrow$ $\Re s=\tfrac12$;
\item \textbf{flips chirality across the line}
      (\leanid{octTorsionWeightGap\_pos\_iff},
      \leanid{octTorsionWeightGap\_neg\_iff}): the defect points along
      $-e_3-e_4$ for $\sigma<\tfrac12$ and along $+e_3+e_4$ for
      $\sigma>\tfrac12$ — the critical line is the unique
      chirality-neutral locus, the CP-odd signature of the channel;
\item \textbf{SO(8) orbit law} (\leanid{octTorsion\_so8\_orbit\_iff}):
      FE partners' torsions lie on a common SO(8) orbit (equal norms)
      iff $\Re s=\tfrac12$ — and since both vectors share one ray, orbit
      coincidence already forces vector equality
      (\leanid{octTorsion\_orbit\_eq\_iff\_collapse}).
\end{itemize}

\paragraph{Zeros on the critical line --- scope of the equivalence.}
\leanid{RH\_iff\_zero\_torsion\_cancellation}: RH $\Leftrightarrow$ at
every nontrivial zero the FE torsion defect vanishes for every triple.
The zeros are exactly the points where the SO(8) twist closes the torsion
loop.  Because the $S^7$ direction is constant, cancellation is governed
by the radial weight alone, and the radial weight is $\sigma$-driven — so
the twist sharpens the geometry (single $S^7$ point, odd defect,
chirality flip, orbit equivalence) without changing the logical status:
``torsion cancels at every zero'' is an exact RH reformulation with no extra
conditions.  (Frozen-table caveat: the witness direction $-e_3-e_4$ is
specific to the current octonion tables; see
\texttt{AGENTS/OCTONION\_TABLE\_AUDIT\_TODO.md}.)

\subsection{The log edge: phases, Diophantine independence, and the
Goldbach circle}\label{sec:log-edge}

Every equivalent condition so far has been a modulus formula.
\leanfile{Hqiv/Story/S3LogPhaseEdge.lean} formalizes why that matters and
what lives on the other side.

\paragraph{Polar decoupling.}
The carrier factorizes exactly
(\leanid{so4SpectralLine\_polar}, \leanid{log\_edge\_decoupling}):
\[
  n^{-s} \;=\; n^{-\sigma}\,\cdot\,e^{-it\log n},
  \qquad
  \lVert e^{-it\log n}\rVert = 1 .
\]
The modulus channel knows only $\sigma$; the unimodular phase channel
knows only $t$.  This is the machine-checked reason the modulus program
produced RH-equivalences and nothing more: \emph{no modulus formula can
import phase data into a $\sigma$-constraint}.  The log edge itself is
$\texttt{linePhase}(pq)=\texttt{linePhase}(p)\cdot\texttt{linePhase}(q)$
(\leanid{linePhase\_mul}) --- multiplication becomes addition in phase.

\paragraph{Diophantine phase independence.}
For distinct primes $p\neq q$, the relation $k\log p=m\log q$ (naturals
or integers) forces $k=m=0$
(\leanid{prime\_log\_nat\_rel}, \leanid{prime\_log\_int\_rel}), proved
from unique factorization ($p^k=q^m$ is impossible).  Consequently
\textbf{two prime phases pin the height}
(\leanid{two\_prime\_phases\_pin\_height}): if the phases of two distinct
primes agree at heights $t_1,t_2$, then $t_1=t_2$.  One prime leaves a
$2\pi/\log p$ lattice of ghost heights; two incommensurable primes leave
none.

\paragraph{The Goldbach circle: additive curvature.}
The additive channel supplies literal curvature
(\leanid{goldbach\_pair\_circle}): a midpoint pair $p+q=2N$ satisfies
\[
  pq + (q-N)^2 \;=\; N^2 ,
\]
so Goldbach pairs are points on the circle of radius $N$ in the
(geometric-mean, deviation) plane.  The multiplicative phase speed is
capped by the additive constraint: $\log(pq)\le2\log N$, with equality
exactly at the diagonal $p=q=N$ (\leanid{pair\_phase\_speed\_max}).

\paragraph{Zeros are addressed by two prime phases.}
Combining the channels
(\leanid{RH\_zero\_determined\_by\_two\_prime\_phases}): under RH, a
nontrivial zero is \emph{completely determined} by the phases of any two
distinct primes at its height --- $\sigma$ from RH, $t$ from phase
independence.  This is the precise form of the ``transformer that locates
a zero'': feed it two prime phases, receive the zero.

\paragraph{The frontier, named.}
The curvature is real, but it lives entirely in the $t$/additive channel,
and decoupling cleanly separates it from $\sigma$.  Any genuine forcing
attack must couple the two channels --- which is what the classical
explicit formula does analytically through the zero set itself, and what
remains open here.  We name that coupling as the frontier rather than
claim it.

\section{Explicit \texorpdfstring{$h(s)$}{h(s)} on the strip}\label{sec:explicit-h}

On the open strip with $\sigma\neq\tfrac12$ the assembly collapses to the
explicit quotient (Lean:
\leanfile{Hqiv/Story/S3InteriorStripHClosedForm.lean}):
\begin{equation}\label{eq:h-explicit}
  h(s)
  =\frac{\zeta(s)}{\texttt{so4CriticalFactor}(s)}
  =\frac{\zeta(s)\,\sqrt2}{\,2\sigma-1\,}.
\end{equation}
Substituting the functional-equation assembly~\eqref{eq:fe-strip} before the
quotient gives the full closed form with $\Gamma(1-s)$ and
$\cos(\pi(1-s)/2)$ in the numerator
(\leanid{interiorStripH\_fe\_closed\_form}): the apparent regional singularities
are then \emph{visible in $h$ itself}, not hidden inside $\zeta$ alone.

\begin{theorem}[Reflection symmetry of the critical factor]
For all $s\in\mathbb{C}$,
\[
  \texttt{so4CriticalFactor}(1-s)=-\texttt{so4CriticalFactor}(s),
\]
because the readout depends only on $\sigma$ and swaps as $\sigma\leftrightarrow1-\sigma$
(\leanid{so4CriticalFactor\_one\_sub}).
\end{theorem}

\begin{proposition}[Reflection product seed (pathway A)]
For $0<\sigma<1$, $\sigma\neq\tfrac12$,
\[
  h(s)\,h(1-s)
  =-\frac{\zeta(s)\,\zeta(1-s)}{\texttt{so4CriticalFactor}(s)^2}.
\]
\end{proposition}

Multiplying $h$ by the equator factor recovers $\zeta$ on the strip
(\leanid{completedInteriorFromH\_eq\_zeta\_on\_strip}); the main conditional predicate
\texttt{InteriorStripHNonvanishingCapstone} is \textbf{logically equivalent} to
\texttt{RiemannHypothesis}
(\leanid{interior\_capstone\_iff\_RiemannHypothesis})---not a weaker statement.

\section{Interior factorization and RH status}\label{sec:rh-status}

On $0<\Re s<1$, $\sigma\neq\tfrac12$, define
\[
  h(s):=\texttt{interiorStripH}(s)=\frac{\zeta(s)}{\texttt{so4CriticalFactor}(s)}.
\]
Then $\zeta(s)=h(s)\cdot(2\sigma-1)/\sqrt2$ off the line
(\leanfile{Hqiv/Story/S3SO4InteriorWitness.lean}).

\begin{definition}[Main RH conditional hypothesis]
\texttt{InteriorAssemblyNonzeroAtNontrivialZerosOffLine} $h$ asserts:
if $\zeta(s)=0$ for a nontrivial zero with $\sigma\neq\tfrac12$, then $h(s)\neq0$.
\end{definition}

\begin{theorem}[Conditional RH]
If this hypothesis holds for \texttt{interiorStripH}, then
\texttt{RiemannHypothesis} follows
(\leanid{RiemannHypothesis\_of\_SO4\_interior\_witness}).
\end{theorem}

\begin{remark}[Honesty]
The factorization identity is proved; this hypothesis is \emph{not}.  Global
factorization on the entire strip including $\sigma=\tfrac12$ is false because
$\zeta(\tfrac12+it)$ is generally nonzero while the critical factor vanishes.
The quotient witness $h=\zeta/\texttt{so4CriticalFactor}$ proves algebra on
off-line points only; it does not by itself prove the hypothesis.
\end{remark}

The even channel \texttt{evenStripChannel} is presently degenerate on the strip
because the odd FE channel already carries the full $\zeta$; the physical even
content is the harmonic--$\Delta$ multiplier $1+\alpha/3$ and the $\pi^4/3$
sphere orbit (\leanfile{Hqiv/Story/S3HarmonicDeltaEvenOrbit.lean}).  A
non-degenerate even/odd split is the remaining analytic packaging target.

\section{Vaporization, backprojection, and the thermodynamic arrow}\label{sec:vaporization}

The de Bruijn--Newman picture (Rodgers--Tao: $\Lambda\ge0$, hence RH
$\Leftrightarrow\Lambda=0$) reads as a heat flow that \emph{vaporizes} off-line
zeros forward in time; RH holds exactly when the present-day function sits at
the vaporization front, i.e.\ there is \textbf{no backprojection margin}
($\Lambda\le0$).  \leanfile{Hqiv/Story/S3HeatFlowArrowNoBackprojection.lean}
pollinates this picture with the already-proved HQIV temperature ladder and
thermodynamics.

\subsection{Dilation slot: where $\tau$ lives in $\mathfrak{so}(4,2)$}

The conformal closure of the SO(4) carrier is
$\mathfrak{so}(4,2)=\mathfrak{so}(4)\oplus D\oplus P_4\oplus K_4$ with dimension
count $15=6+4+4+1$ (\leanid{conformal\_dilation\_slot\_count}); the chain
$\mathfrak{so}(4)\subset\mathfrak{so}(4,2)\subset\mathfrak{so}(8)$ is
$6\le15\le28$ (\leanid{so4\_so42\_so8\_dim\_chain}).  The heat-flow deformation
time $\tau$ occupies the \emph{one-dimensional dilation slot} $D$;
backprojection is negative dilation.  This is the same carrier escalation that
produced the $45^\circ$ equator readout, now read in the time channel.

\subsection{Backprojection dichotomy (proved)}

On the discrete ladder weight
$w_\tau(m)=\exp(-\tau\,t_{\mathrm{HQIV}}(m))$:
\begin{itemize}[nosep]
\item \textbf{forward} ($\tau\ge0$): $w_\tau(m)\le1$ on every shell and the
      deformed shell sum stays summable on $\Re s>1$
      (\leanfile{Hqiv/Physics/HQIVHeatFlowDeformation.lean});
\item \textbf{backward} ($\tau<0$): $w_\tau(m)>1$ off the horizon shell
      (\leanid{backprojection\_weight\_gt\_one}), grows strictly shell by shell
      (\leanid{backprojection\_weight\_strict\_growth}), and is
      \emph{unbounded} along the ladder
      (\leanid{backprojection\_weight\_unbounded}): no constant dominates a
      backward semigroup on the discrete carrier.
\end{itemize}
The proven second law (entropy-production positivity,
\leanfile{Hqiv/Physics/ThermodynamicLawsFromLadder.lean}) makes the
arrow-induced deformation parameter nonnegative, so the physical flow can only
move \emph{toward} vaporization
(\leanid{arrow\_forbids\_backprojection},
\leanid{arrow\_weight\_contractive}).

\subsection{$\lambda$-lock reduction and the honest frontier}

The ladder bundle carries $\lambda_{\mathrm{HQIV}}\ge0$ --- the analogue of the
Rodgers--Tao inequality.  Exactly as $\Lambda\ge0$ reduces RH from
$\Lambda=0$ to the one-sided bound $\Lambda\le0$, the analogue zero-lock is
proved equivalent to the no-backprojection inequality:
\[
  \lambda_{\mathrm{HQIV}}=0
  \quad\Longleftrightarrow\quad
  \lambda_{\mathrm{HQIV}}\le0
\]
(\leanid{lambdaHQIV\_zero\_iff\_no\_backprojection}).  Concrete finite-window
ladder witnesses already discharge $\lambda_{\mathrm{HQIV}}=0$
(\leanid{lambdaHQIV\_eq\_zero\_of\_finiteWindowConcrete}).

The genuine analytic frontier is the \emph{identification} of the discrete
ladder flow with the classical $\xi$ heat flow.  It is named as the payload
\texttt{VaporizationForcesCriticalLine} and proved \textbf{equivalent to RH}
(\leanid{vaporization\_iff\_RiemannHypothesis}) --- the same honesty pattern as
the Weil-positivity slot.  Bundling the proved $\lambda$-lock side with the
payload gives the \texttt{HeatFlowVaporizationBridge}; given any concrete
ladder witness,
\[
  \texttt{HeatFlowVaporizationBridge}\ \text{inhabited}
  \quad\Longleftrightarrow\quad
  \texttt{RiemannHypothesis}
\]
(\leanid{vaporizationBridge\_iff\_RiemannHypothesis}): the thermodynamic side is
free, and the vaporization payload is the entire remaining content.

\subsection{Classical target formalized: $\Phi$, $H_t$, and lattice inheritance}

\leanfile{Hqiv/Story/S3DBNLatticeInheritance.lean} makes the classical side
concrete in Lean, so the inheritance claim is a precise statement rather than
prose:
\begin{itemize}[nosep]
\item \textbf{The density.}  The de Bruijn--Newman density
      $\Phi(u)=\sum_{m\ge1}(2\pi^2m^4e^{9u}-3\pi m^2e^{5u})e^{-\pi m^2e^{4u}}$
      is \texttt{dbnPhi}; for $u\ge0$ every term is \emph{proved} nonnegative
      (\leanid{dbnPhiTerm\_nonneg}; the quartic shell dominates as soon as
      $2\pi m^2e^{4u}\ge3$, which holds on every shell since $\pi>3$) and the
      series is \emph{proved} summable
      (\leanid{dbnPhi\_summable}; polynomial-times-geometric domination with
      ratio $e^{-\pi e^{4u}}<1$).
\item \textbf{The deformed family.}
      $H_t(z)=\int_{u>0}e^{tu^2}\,\Phi(u)\cos(zu)\,du$ is
      \texttt{dbnHeatFamily}; $H_0$ is the (normalized) $\xi$.
\item \textbf{Weight-level dictionary (proved).}  The classical Gaussian
      factor \emph{is} the repo's discrete heat weight under
      $(\tau,u)\mapsto(-t,u^2)$:
      \texttt{dbnGaussianFactor}$\,t\,u=$
      \texttt{discreteHeatKernelWeight}$\,(-t)\,(u^2)$
      (\leanid{dbnGaussianFactor\_eq\_discreteHeatKernelWeight}).  The
      dichotomy transfers: $t\le0$ contracts every mode
      (\leanid{dbn\_backward\_factor\_le\_one}), $t>0$ amplifies without bound
      (\leanid{dbn\_forward\_factor\_unbounded}) --- the continuum mirror of
      the lattice backprojection blow-up.
\end{itemize}
The bridge \texttt{DBNLatticeInheritance} then carries the classical constant
$\Lambda$ with the two classical theorems as explicit fields (Rodgers--Tao
$\Lambda\ge0$; Newman/de Bruijn RH $\Leftrightarrow\Lambda\le0$ --- neither yet
in Mathlib) plus the single new claim
$\Lambda\le\lambda_{\mathrm{HQIV}}$.  Since the ladder witness proves
$\lambda_{\mathrm{HQIV}}=0$, the bridge forces $\Lambda=0$
(\leanid{Lambda\_eq\_zero\_of\_inheritance}) and hence RH
(\leanid{RiemannHypothesis\_of\_dbnLatticeInheritance}); it also instantiates
the vaporization bridge
(\leanid{vaporizationBridge\_of\_dbnLatticeInheritance}).

\begin{remark}[The conviction, faithfully named]
Given the two classical imports, the lattice-domination bound is proved
\emph{equivalent} to RH:
\[
  \Lambda\le\lambda_{\mathrm{HQIV}}(W)
  \quad\Longleftrightarrow\quad
  \texttt{RiemannHypothesis}
\]
(\leanid{lattice\_dominates\_iff\_RiemannHypothesis}).  The thermodynamic side
is free and the classical constants are pinned; the entire content of ``the
discrete no-backprojection lock transfers to the continuum'' is the RH
frontier itself --- now a single named inequality between two Lean terms
rather than an informal analogy.
\end{remark}

\subsection{Discrete dominates continuum: comparison inequality and exact
phase lock}\label{sec:dbn-comparison}

\leanfile{Hqiv/Story/S3DBNDiscreteContinuumComparison.lean} pushes the proved
boundary flush against the inheritance frontier with a genuine one-sided
comparison between the two heat families, on the region where both converge
absolutely.  All of the following is unconditional (no hypotheses beyond
parameter ranges).

\paragraph{Analytic control of the classical side.}
The dBN density is \emph{strictly} positive for $u\ge0$
(\leanid{dbnPhi\_pos}), decays super-exponentially,
$\Phi(u)\le C\,e^{-3u}$ via $e^{-\pi m^2e^{4u}}\le
e^{-\pi(m^2-1)}e^{-\pi e^{4u}}$ and $\pi>3$
(\leanid{dbnPhi\_le\_decay}), is continuous on $u\ge0$ by a Weierstrass
M-test (\leanid{dbnPhi\_continuousOn}), and is Lebesgue-integrable on
$(0,\infty)$ (\leanid{dbnPhi\_integrableOn}).  Consequently the
\emph{backward} family is uniformly bounded on the real axis:
\[
  t\le0,\ z\in\mathbb{R}
  \quad\Longrightarrow\quad
  \lVert H_t(z)\rVert \;\le\; \int_{u>0}\Phi(u)\,du
\]
(\leanid{dbnHeatFamily\_norm\_le}), while the forward family is unbounded
(\leanid{dbn\_forward\_factor\_unbounded}).  At the anchor $z=0$ the backward
family is a strictly positive real
(\leanid{dbnWeightedIntegral\_pos}), so its phase is exactly $0$
(\leanid{dbnHeatFamily\_arg\_at\_zero}).

\paragraph{The heat flow fixes the horizon shell.}
Since $t_{\mathrm{HQIV}}(0)=0$, the $m=0$ weight of the discrete deformed sum
is $e^{-\tau\cdot0}=1$ for \emph{every} deformation time $\tau$
(\leanid{hqivHeatKernelWeight\_horizon}).  At trivial twist ($\varphi=0$) and
real anchor $s=\sigma>1$ the full complex sum is the coercion of a convergent
positive real series (\leanid{HQIVDeformedSum\_real\_anchor}), so the
flow-invariant horizon term gives the unconditional lower bound
$\mathrm{eff}_0^{-\sigma}\le\lVert\texttt{HQIVDeformedSum}\rVert$
(\leanid{HQIVDeformedSum\_anchor\_lower}): the flow can never dim the horizon,
and every other shell only adds positive mass.

\paragraph{The comparison inequality and phase alignment.}
Combining the two sides, for every $\tau\ge0$, $t\le0$, $\sigma>1$ and every
real $z$ there is an \emph{explicit} constant
$c=\mathrm{eff}_0^{-\sigma}/(\int\Phi+1)>0$ with
\[
  c\,\lVert H_t(z)\rVert \;\le\;
  \lVert\texttt{HQIVDeformedSum}\,\tau\,\dots\,\sigma\rVert
\]
(\leanid{discrete\_dominates\_continuum\_anchor}): no backward continuum
deformation can outgrow the lattice sum --- the growth sign is preserved.
The phase-alignment statement holds in its strongest form: the arguments of
the discrete sum and the backward continuum family \emph{coincide} at the
anchor, so the cosine of their difference is exactly $1$
(\leanid{phase\_alignment\_anchor}), in particular bounded below by $\tfrac12$
(\leanid{phase\_alignment\_cos\_ge\_half}).

\begin{remark}[Honest scope]
The comparison lives on the absolutely-convergent region ($\Re s>1$ discrete
$\times$ real axis continuum, all $t\le0$).  Extending it into the critical
strip \emph{with zero tracking} is precisely the inheritance bound
$\Lambda\le\lambda_{\mathrm{HQIV}}$, which
\leanid{lattice\_dominates\_iff\_RiemannHypothesis} shows is RH itself.  The
module advances the proved boundary up to that frontier; it does not cross it.
\end{remark}

\subsection{The constant itself: $\Lambda$ ported as a Lean definition}
\label{sec:dbn-port}

\leanfile{Hqiv/Story/S3DBNConstantPort.lean} removes the last abstraction from
the inheritance bridge: the classical constant is no longer an arbitrary real
packaged in a structure field but a \emph{concrete Lean term},
\[
  \texttt{dbnLambda} \;=\;
  \inf\,\{\,t\in\mathbb{R} \mid \forall z,\ H_t(z)=0\Rightarrow\Im z=0\,\},
\]
the infimum of the hyperbolicity locus of the very \texttt{dbnHeatFamily}
whose Gaussian factor was proved identical to the lattice heat kernel.  The
two classical imports are now precisely typed statements about \emph{that}
constant --- the form required for an actual Mathlib contribution:
\begin{itemize}[nosep]
\item \texttt{RodgersTaoStatement} $:= 0\le\texttt{dbnLambda}$
      (Rodgers--Tao, Annals 191 (2020));
\item \texttt{NewmanDBNStatement} $:=$ \texttt{RiemannHypothesis}
      $\Leftrightarrow \texttt{dbnLambda}\le0$ (Newman 1976; de Bruijn 1950);
\item \texttt{DBNUpwardClosure} --- de Bruijn's heat-flow preservation of
      hyperbolicity (the locus is an upper ray): the continuum analytic
      mirror of the lattice no-backprojection dichotomy;
\item \texttt{DBNXiBridge} --- the normalization
      $H_0(z)\propto\xi(\tfrac12+\tfrac{iz}2)$ against Mathlib's
      \texttt{completedRiemannZeta}, stated normalization-robustly.
\end{itemize}
Unconditionally proved on top of the definition: the backward family is
nondegenerate at the anchor, $H_t(0)\ne0$ for $t\le0$
(\leanid{dbnHeatFamily\_ne\_zero\_at\_anchor}); the $\inf$ interface
(\leanid{dbnLambda\_le\_of\_realZerosAt},
\leanid{le\_dbnLambda\_of\_forall}); upper-set propagation under de Bruijn
closure (\leanid{dbnRealZerosAt\_mono\_of\_upwardClosure}); and, given the two
classical statements, $\texttt{dbnLambda}=0\Leftrightarrow$ RH
(\leanid{dbnLambda\_eq\_zero\_iff\_RH}).  The honesty theorems of the
inheritance bridge now hold verbatim at the concrete constant:
\[
  \texttt{dbnLambda}\le\lambda_{\mathrm{HQIV}}(W)
  \quad\Longleftrightarrow\quad
  \texttt{RiemannHypothesis}
\]
(\leanid{concrete\_lattice\_dominates\_iff\_RiemannHypothesis},
\leanid{concrete\_dbnLatticeInheritance\_iff\_RiemannHypothesis}).  Porting
the \emph{proofs} of Rodgers--Tao and Newman/de Bruijn (Laguerre--P\'olya
theory, heat-flow zero dynamics) remains a major formalization programme; the
objects and statements now have their precise Lean targets.

\section{The shared $\tfrac12$ bridge: Riemann Hypothesis and Goldbach}\label{sec:rh-goldbach}

\subsection{Two classical problems, one projection constant}

The Riemann Hypothesis and Goldbach's conjecture (even parity case: every even
$n>2$ is a sum of two primes) are unrelated in classical analytic number theory.
In the HQIV $\mathfrak{so}(8)$ / SO(4) projection story they meet on the same
normalized constant $\tfrac12$:
\begin{center}
\begin{tabular}{@{}lll@{}}
\toprule
Channel & Geometric object & $\tfrac12$ meaning \\
\midrule
$\zeta$ / zeros & critical line $\Re s=\tfrac12$ &
  \texttt{rot45Free}$(\sigma)=0\Leftrightarrow\sigma=\tfrac12$ \\
Prime sums & midpoint $N$, pair $(p,q)$ with $p+q=2N$ &
  tangent slope $N/(p+q)=\tfrac12$ \\
\bottomrule
\end{tabular}
\end{center}
The $\zeta$ row is the pointwise equator condition of \S\ref{sec:45-proj}.  The
Goldbach row is proved \emph{for any certified midpoint pair}, without assuming
Goldbach:
\[
  \frac{N}{p+q}=\frac{N}{2N}=\frac12
  \quad\text{whenever }p+q=2N
\]
(\leanfile{Hqiv/Geometry/GoldbachG2Parity.lean},
\leanid{so4\_orthogonal\_tangent\_midpoint\_slope\_eq\_half}).
Diagonal partitions $N+N=2N$ when $N$ is prime carry the same slope; there is no
exception branch (\leanid{so4\_orthogonal\_diagonal\_tangent\_slope\_eq\_half}).

\subsection{Midpoint Goldbach and the holonomy bridge}

Lean names the arithmetic target in \emph{midpoint} form: for each midpoint $N$,
primes $p\le N\le q$ with $p+q=2N$ (\texttt{GoldbachMidpointPair}).  Eventual
midpoint Goldbach above a threshold $N_0$ is
\texttt{MidpointGoldbachEventually}~$N_0$; even Goldbach parity
(\texttt{GoldbachParity}) follows at $N_0=2$ via
\texttt{goldbach\_pair\_of\_midpoint\_pair}.

The open arithmetic--geometric bridge is positivity of the finite paired-prime
count around each midpoint:
\[
  \texttt{goldbachMidpointCount}(N)>0
  \quad\Longrightarrow\quad
  \exists\,p,q:\ \texttt{GoldbachMidpointPair}\,N\,p\,q,
\]
packaged as \texttt{SO4ZetaHolonomyForcesMidpointPairs}~$N_0$
(\leanid{midpoint\_goldbach\_of\_so4\_zeta\_holonomy\_bridge}).
Any certified midpoint pair lifts to integrable Hopf-shell holonomy with
$\Delta$/G$_2$/triality fields
(\leanid{hopf\_fiber\_midpoint\_holonomy\_support\_of\_midpoint\_pair}).

The constructive $G_2+\Delta$ harmonic landing layer
(\leanfile{Hqiv/Geometry/GoldbachG2Parity.lean}) states the complementary
finite certificate target \texttt{DeltaHarmonicCompleteness}: once every even
$n>2$ has a $\Delta$-harmonic landing, \texttt{goldbach\_from\_delta\_harmonic}
extracts \texttt{GoldbachParity} with no extra arithmetic input.

\subsection{The bridge \emph{is} RH $\wedge$ Goldbach: proved equivalence}

\leanfile{Hqiv/Story/S3ExplicitFormulaDualitySlot.lean} names the shared payload
explicitly.  An \texttt{SO8ProjectedHalfSlopeBridge}~$N_0$ bundles:
\begin{itemize}[nosep]
\item \textbf{$\zeta$ side.} \texttt{WeilPositivityForcesCriticalLine}, i.e.\ the
      explicit-formula / positivity localization step equivalent to RH
      (\leanid{weilPositivity\_iff\_RiemannHypothesis});
\item \textbf{prime side.} \texttt{SO4ZetaHolonomyForcesMidpointPairs}~$N_0$.
\end{itemize}
Both directions are now machine-checked, so the packaging is a proved
\emph{logical equivalence}, not merely a sufficient condition:
\begin{align}
  \texttt{SO8ProjectedHalfSlopeBridge}\,N_0
  &\Longleftrightarrow
  \texttt{RiemannHypothesis}
  \wedge
  \texttt{MidpointGoldbachEventually}\,N_0,
  \label{eq:bridge-iff-general}\\
  \texttt{SO8ProjectedHalfSlopeBridge}\,2
  &\Longleftrightarrow
  \texttt{RiemannHypothesis}
  \wedge
  \texttt{GoldbachParity}.
  \label{eq:bridge-iff-two}
\end{align}
Equation~\eqref{eq:bridge-iff-general} is
\leanid{so8\_projected\_half\_slope\_iff\_rh\_and\_midpoint\_goldbach};
equation~\eqref{eq:bridge-iff-two} is the capstone
\leanid{so8\_projected\_half\_slope\_two\_iff\_rh\_and\_goldbach\_parity}.
The converse direction is built from two reconstruction lemmas:
\begin{itemize}[nosep]
\item a certified midpoint pair forces a positive midpoint count
      (\leanid{midpoint\_count\_pos\_of\_midpoint\_pair}), so the holonomy
      payload is \emph{equivalent} to eventual midpoint Goldbach
      (\leanid{so4\_zeta\_holonomy\_bridge\_iff\_midpoint\_goldbach});
\item an unordered Goldbach pair for $2N$ reorders into a midpoint pair at $N$
      (\leanid{midpoint\_pair\_of\_goldbach\_pair\_two\_mul}), giving
      \texttt{MidpointGoldbachEventually}~$2\Leftrightarrow
      \texttt{GoldbachParity}$
      (\leanid{midpoint\_goldbach\_two\_iff\_goldbach\_parity}).
\end{itemize}
The channels are also separated: the $\zeta$ field alone is exactly RH, and the
midpoint field at threshold~$2$ alone is exactly Goldbach parity, with no
leakage between them
(\leanid{bridge\_channels\_are\_rh\_and\_goldbach},
\leanid{so4\_zeta\_holonomy\_bridge\_two\_iff\_goldbach\_parity}).

\begin{remark}[What ``RH $=$ Goldbach'' means here]
This note does \emph{not} claim a classical theorem $\text{RH}\Leftrightarrow
\text{Goldbach}$ in Mathlib.  What \emph{is} proved is the identification of
the geometric payload with the conjunction: the half-slope bridge is
\textbf{provably the same proposition} as
$\text{RH}\wedge\text{Goldbach parity}$
(equation~\eqref{eq:bridge-iff-two}), with the same normalized $\tfrac12$
readout serving as the critical-line lock on the $\zeta$ channel and the
midpoint tangent slope on the prime channel.  The packaging therefore has zero
slack: discharging the bridge is exactly as hard as the two open problems
combined---by the equivalence theorem, \emph{provably} so---and conversely any
future proof of RH and Goldbach immediately instantiates the bridge.  Neither
classical conjecture is proved here; what is discharged is the faithfulness of
the geometric encoding itself.
\end{remark}

\subsection{Prime side of the explicit formula}

The von Mangoldt weight $\Lambda(n)=\log p$ on prime powers is the concrete
prime-power content dual to the zero sum in the explicit formula
(\leanfile{Hqiv/Story/S3ExplicitFormulaDualitySlot.lean},
\texttt{chebyshevPsi}, \leanid{vonMangoldt\_prime\_eq\_log}).  Positivity on
that prime side is the standard route to critical-line zero localization; in this
repo it is named faithfully as RH-equivalent, parallel to the interior capstone
on $h(s)$.

\section{PTE degree~3 and SO(4) cancellation}\label{sec:pte-parallel}

Recent work of Lee, Takahashi, and Tsai~\cite{LeeTakahashiTsai2026PTE} shows that
anomaly cancellation in a chiral hidden sector with unbroken vector-like
$U(1)_{\mathrm H}\times U(1)_{\mathrm X}$ is \emph{equivalent} to the degree
$k=3$ Prouhet--Tarry--Escott (PTE) problem: two equal-size multisets
$A=\{a_i\}$, $B=\{b_i\}$ must satisfy
\[
  \sum_i a_i^{\ell}=\sum_i b_i^{\ell}
  \qquad(\ell=1,2,3).
\]
Minimal ideal solutions have size $n=k+1=4$, matching the four-dimensional
SO(4) quaternion carrier of \S\ref{sec:harmonic-delta}.  The phenomenological
prediction---at least four hidden mass eigenstates, typically organized as
near-degenerate doublets when the PTE solution is sign-symmetric---parallels
the six-pole geometry of \S\ref{sec:quaternion}: antipodal pairs
$(+i,-i)$, $(+j,-j)$, $(+k,-k)$ cancel as \emph{orbits} while each pole retains
a nonzero $\pm$ residual (\leanfile{Hqiv/Story/S3SixPolesResidual.lean}).

\subsection{Sign-symmetric multisets and head/tail reflection}

Lee \emph{et al.}\ focus on \textbf{symmetric} PTE solutions
$A=\{\pm a_1,\ldots,\pm a_m\}$, $B=\{\pm b_1,\ldots,\pm b_m\}$.
Under this sign pairing, the $\ell=1$ and $\ell=3$ constraints are automatic
(both sides sum to zero), so the only nontrivial ideal-$n=4$ condition is
\[
  a_1^2+a_2^2=b_1^2+b_2^2,
\]
a matching of \emph{even} (quadratic) invariants.
This is the same parity split as head/tail reflection on the S$^3$ shell:
\[
  \texttt{criticalProj}(p)+\texttt{criticalProj}(\texttt{headTailReflect}\,p)=0
\]
(\leanid{headTail\_orbit\_pair\_cancels}) forces \textbf{linear} ($\ell=1$)
cancellation for every orbit pair, while single-axis poles never vanish
pointwise (\leanid{criticalProj\_ne\_zero\_of\_singleAxis}).
Geometrically, the PTE even constraint is orthogonal invariance of a
two-vector norm; the $45^\circ$ readout in \S\ref{sec:45-proj} is the same
rotation that fixes the diagonal slot at $1/\sqrt2$ and isolates the equator
slot $(2\sigma-1)/\sqrt2$.

\subsection{Orbit vs.\ pointwise: two readings of ``cancellation''}

Lee \emph{et al.}\ distinguish PTE solutions related by affine charge shifts
(number-theoretically equivalent, physically distinct mass/coupling assignments)
from the orbit-level doublet structure their symmetric solutions predict.
The S$^3$ story makes the same distinction internally:
\begin{itemize}[nosep]
\item \textbf{Orbit cancellation} --- reflection pairs $\{\sigma,1-\sigma\}$ and
      head/tail pairs always sum to zero in the free coordinate
      (\leanfile{Hqiv/Story/S3OrbitVsPointwiseGap.lean},
      \leanid{orbit\_free\_sum\_cancels}); this holds even off
      $\Re s=\tfrac12$.
\item \textbf{Pointwise cancellation} --- a \emph{single} slot vanishes iff
      $\sigma=\tfrac12$ (Theorem~\ref{thm:critical-line}); upgrading orbit
      cancellation to pointwise cancellation at actual $\zeta$ zeros is the RH
      capstone (\S\ref{sec:rh-status}).
\end{itemize}
Reading PTE through this lens: matching power sums is the discrete
charge-theoretic avatar of orbit cancellation; selecting a \emph{physical}
spectrum (mass doublets, minicharge partners) is the pointwise upgrade---an
analogue of RH-hard localization on the critical line.

\subsection{Extension toward hidden-sector model building}

Three concrete bridges from this note to~\cite{LeeTakahashiTsai2026PTE}:
\begin{enumerate}[label=(\roman*),nosep]
\item \textbf{Lean packaging.}  Formalize symmetric PTE-$k=3$ solutions as
      charge multisets with proved automatic $\ell=1,3$ cancellation under sign
      pairing, parallel to \leanfile{Hqiv/Story/S3PrimeAxisCancellation.lean}.
\item \textbf{Ideal $n=4$ sum-of-squares.}  Embed minimal PTE solutions as
      integer samples on the $\pm i,\pm j$ (or $\pm j,\pm k$) pole lattice;
      the surviving constraint $a_1^2+a_2^2=b_1^2+b_2^2$ is then a
      two-axis norm match on $S^3$.
\item \textbf{TUFT hidden readout.}  In the HQIV patch ontology, finite SM
      anomaly traces are proved at the coefficient level
      (\leanfile{Hqiv/Algebra/AnomalyCancellation.lean}); extending the same
      finite-sum discipline to $U(1)_{\mathrm H}\times U(1)_{\mathrm X}$ content
      would import Lee \emph{et al.}'s PTE equivalence as a \emph{model-building
      corollary} of the SO(4) cancellation geometry, not as an external axiom.
\end{enumerate}

\section{Pathways toward an unconditional proof}\label{sec:pathways}

The analysis below matches the claim-status table: the capstone is RH-hard.
Plausible extensions that exploit \emph{this} framework (not imported from
external regularization claims) include:
\begin{enumerate}[label=\textbf{\Alph*.},nosep]
\item \textbf{Completed $h$.}  Use~\eqref{eq:h-explicit} and
      \texttt{so4CriticalFactor}$(1-s)=-$ \texttt{so4CriticalFactor}$(s)$ to
      build a candidate entire/completed interior function; check whether
      reflection forces off-line zeros to contradict discharge lemmas.
\item \textbf{Zero-free re-expression.}  Rewrite $h$ as a manifestly nonzero
      Dirichlet/Euler or harmonic-weighted integral on $\sigma\neq\tfrac12$.
\item \textbf{SO(4)/patch obstruction.}  Relate an off-line zero to a nonzero
      holonomy or commutator on the octonion carrier, contradicting proved
      $\mathfrak{so}(8)$ closure.
\item \textbf{Classical tools in channel language.}  Phragm\'en--Lindel\"of on
      vertical lines $\sigma\neq\tfrac12$ using the $\tan(\pi s/2)$ ratio;
      explicit-formula with separated sin/cos slots.
\item \textbf{Non-degenerate even split.}  Promote the harmonic--$\Delta$
      multiplier into \texttt{evenStripChannel} so $h$ is a genuine even/odd
      assembly rather than a pure quotient.
\item \textbf{Shared half-slope bridge (\S\ref{sec:rh-goldbach}).}  Prove
      \texttt{SO4ZetaHolonomyForcesMidpointPairs} from finite patch holonomy on
      the same $\tfrac12$ readout that locks $\zeta$-zeros to $\Re s=\tfrac12$,
      or prove \texttt{DeltaHarmonicCompleteness} on the $G_2+\Delta$ landing
      channel.
\item \textbf{PTE hidden-sector packaging (\S\ref{sec:pte-parallel}).}
      Import~\cite{LeeTakahashiTsai2026PTE} symmetric PTE-$k=3$ solutions as
      charge samples on the six-pole lattice; use sum-of-squares matching as the
      even-channel counterpart to $\texttt{rot45Diag}$ rigidity.
\end{enumerate}

\section{Lean module index}\label{sec:lean-index}

Build target: \texttt{lake build HQIVStory}.  Core modules (import order in
\leanfile{HQIVStory.lean}):
\begin{itemize}[nosep]
\item \leanfile{Hqiv/Story/S3ClosureDeltaLiftBridge.lean} --- harmonic $H_n\le K(n)$, SO(4) lift;
\item \leanfile{Hqiv/Story/S3HarmonicPrimeZetaPath.lean} --- harmonic $\Rightarrow$ primes $\Rightarrow\zeta\Rightarrow\Lambda$;
\item \leanfile{Hqiv/Story/S3FortyFiveProjection.lean} --- $45^\circ$ diagonal / equator;
\item \leanfile{Hqiv/Story/S3RotationRigidity.lean} --- projection line, twiddle rigidity;
\item \leanfile{Hqiv/Story/S3MirroredTwiddlePair.lean} --- mirrored twiddle pairs
      at $\theta$ and $\pi/2-\theta$, symmetric zero loci, FE orbit cancellation,
      critical-line coincidence;
\item \leanfile{Hqiv/Story/S3CubicPlasticTwiddle.lean} --- higher-twiddle
      factorization probes: linear $\text{free}_c$, plastic
      \texttt{so4PlasticCubicFree}, mirror pairs; not a zero-locator channel
      (\S\ref{sec:higher-twiddle-probes});
\item \leanfile{Hqiv/Story/S3HigherTwiddleFactorizationProbe.lean} ---
      per-channel \texttt{twiddleInteriorAssembly}, patch factorization off probe
      loci, semiprime-supported finite cross-talk, frame-radius obstruction;
\item \leanfile{Hqiv/Story/S3ZetaClosedForm.lean} --- regional closed forms, factorization predicate;
\item \leanfile{Hqiv/Story/S3SO4ZetaProjectionClosedForm.lean} --- $\pi^4/6$ sphere slot, master packaging;
\item \leanfile{Hqiv/Story/S3HarmonicDeltaEvenOrbit.lean} --- $1+\alpha/3=6/5$ orbit;
\item \leanfile{Hqiv/Story/S3ZetaAxisRotationProjection.lean} --- j/k cancellation and ratios;
\item \leanfile{Hqiv/Story/S3SO4InteriorWitness.lean} --- even/odd channels, \texttt{interiorStripH};
\item \leanfile{Hqiv/Story/S3FESlotDischarge.lean} --- negative-integer slot discharge;
\item \leanfile{Hqiv/Story/S3DeltaOrbitOffStrip.lean} --- off-strip confinement;
\item \leanfile{Hqiv/Story/S3SigmaReadoutScope.lean} --- $\sigma$-blindness, no global
      $\zeta=$ readout;
\item \leanfile{Hqiv/Story/S3ZeroQuadrupletOrbit.lean} --- FE pairs, Schwarz quadruplet;
\item \leanfile{Hqiv/Story/S3OrbitVsPointwiseGap.lean} --- orbit vs pointwise RH gap;
\item \leanfile{Hqiv/Story/S3RHDischarge.lean} --- discharge packaging;
\item \leanfile{Hqiv/Story/S3InteriorStripHClosedForm.lean} --- explicit $h(s)$,
      reflection symmetry, capstone $\Leftrightarrow$ RH;
\item \leanfile{Hqiv/Story/S3ExplicitFormulaDualitySlot.lean} --- prime side
      $\psi(x)$, \texttt{SO8ProjectedHalfSlopeBridge}, RH $\wedge$ Goldbach joint
      packaging;
\item \leanfile{Hqiv/Geometry/GoldbachG2Parity.lean} --- midpoint slope $\tfrac12$,
      holonomy bridge, $G_2+\Delta$ landing certificates;
\item \leanfile{Hqiv/Story/S3HeatFlowArrowNoBackprojection.lean} --- dilation slot
      count, backprojection dichotomy, arrow $\Rightarrow$ no backprojection,
      $\lambda$-lock, \texttt{HeatFlowVaporizationBridge} $\Leftrightarrow$ RH;
\item \leanfile{Hqiv/Story/S3DBNLatticeInheritance.lean} --- concrete $\Phi$ and
      $H_t$, term positivity and summability, weight-level lattice dictionary,
      \texttt{DBNLatticeInheritance}, inheritance bound $\Leftrightarrow$ RH;
\item \leanfile{Hqiv/Story/S3DBNDiscreteContinuumComparison.lean} --- strict
      positivity, decay, continuity and integrability of $\Phi$; uniform bound
      on backward $H_t$; flow-invariant horizon term; comparison inequality
      $c\lVert H_t(z)\rVert\le\lVert\texttt{HQIVDeformedSum}\rVert$ and exact
      phase alignment;
\item \leanfile{Hqiv/Story/S3DBNConstantPort.lean} --- the de Bruijn--Newman
      constant \texttt{dbnLambda} as a concrete $\inf$, Rodgers--Tao and
      Newman/dBN as precisely typed statements, honesty theorems at the
      concrete constant;
\item \leanfile{Hqiv/Story/S3TangentOrbitCriticalLine.lean} --- tangent law of
      the quadruplet orbit: FE inverts, Schwarz conjugates, modulus identity
      $\lVert\sin\rVert^2-\lVert\cos\rVert^2=-\cos(\pi\sigma)$, critical line
      $=$ unimodular locus, orbit collapse $\Leftrightarrow$ line, equator
      factor tie;
\item \leanfile{Hqiv/Story/S3EulerSpectralCancellation.lean} --- zeros as
      spectral cancellations: nonvanishing Euler factors and no cancellation
      in the product region, breakdown-locus confinement, tame von Mangoldt
      spectrum, resonance blow-up of $\lVert1/\zeta\rVert$ at every zero,
      spectral-cancellation profile with the tangent orbit attached;
\item \leanfile{Hqiv/Story/S3SpectralResonanceChanneling.lean} --- collective
      cancellation (single prime lines and finite truncations never vanish
      for $\Re s>0$), two-channel dichotomy at every zero, geometric-channel
      absorption $\Leftrightarrow$ RH, SO(4) spectral-line transformer with
      FE pairing, geometric-mean and square-root-weight locators, prime
      generation of lines, RH $\Leftrightarrow$ square-root spectral
      weights;
\item \leanfile{Hqiv/Story/S3VerticalAxisRigidity.lean} --- vertical mirror
      and its fixed line, FE closure of the zero set, two-mirrors-make-a-%
      translation rigidity (abstract and zeta-instantiated unique axis),
      FE-partner collinearity pin, square-root-weight locus as a globally
      straight vertical line with flow invariance;
\item \leanfile{Hqiv/Story/S3SameHeightOrbitCollapse.lean} ---
      $\sigma$-uniform collectivity of spectral weights, same-height
      quadruplet partner and
      merge law, RH $\Leftrightarrow$ height pins $\Re$ (FE only), RH
      $\Leftrightarrow$ one zero per height (with Schwarz hypothesis);
\item \leanfile{Hqiv/Story/S3PolarProjectionCollapse.lean} --- polar partner
      and same-height law, polar separation $=$ equator factor
      ($-\sqrt2\cdot\texttt{so4CriticalFactor}$), separation norm $=$ twice
      the line deviation, collapse law, single-point polar projection
      $\Leftrightarrow$ RH, chain to one-zero-per-height;
\item \leanfile{Hqiv/Story/S3ZeroHolonomyGoldbachChain.lean} --- pair
      activation at every zero, AM--GM joint-weight cap $1/(pq)\ge1/N^2$,
      zero-contains-pair-holonomy capstone, parity gives holonomy support,
      bridge $\Leftrightarrow$ spectral weights $\wedge$ parity
      $\Leftrightarrow$ polar collapse $\wedge$ parity;
\item \leanfile{Hqiv/Story/S3TuftNestedFrameTower.lean} --- level-$N$
      spectral frame radius, nesting recursion, $\sigma$-blind flow
      rigidity, positivity, harmonic radius law at every level $N\ge2$,
      RH $\Leftrightarrow$ harmonic tower radii at every zero, divergence
      of the ladder on the line, curvature domination $H_N\le K(N)$;
\item \leanfile{Hqiv/Story/S3DiophantineTransformer.lean} --- level-$N$
      diagonal transformer $D_N(s)$, trace ladder $\to\zeta$ on
      $\Re s>1$, HS norm $=$ frame radius, FE composition $=$ fixed
      rational harmonic operator $\mathrm{diag}(1/n)$ with trace $H_N$,
      adjoint law $D_N(1-s)=D_N(s)^{\dagger}\Leftrightarrow$ line,
      scaled isometry on the line, RH $\Leftrightarrow$ adjoint law at
      every zero;
\item \leanfile{Hqiv/Story/S3OctonionicAssociatorChannel.lean} ---
      commutative saturation of the transformer family (abelian, normal),
      trilinear associator scaling, concrete witness
      $[e_1,e_2,e_4]=-e_3-e_4$ with norm$^2=2$, associator channel
      $2(abc)^{-2\sigma}$, positivity at every zero, square-root locator,
      FE pairing $4/(abc)^2$, RH $\Leftrightarrow$ torsion weights at
      zeros, Goldbach floor $1/N^3$ on the line;
\item \leanfile{Hqiv/Story/S3OctonionS7TorsionCancellation.lean} ---
      torsion vector $=(abc)^{-\sigma}\cdot(-e_3-e_4)$ (single S$^7$
      point), FE defect closed form and oddness, cancellation
      $\Leftrightarrow$ line, chirality flip across the line, SO(8)
      orbit law and orbit $\Rightarrow$ vector collapse, RH
      $\Leftrightarrow$ torsion cancellation at every zero;
\item \leanfile{Hqiv/Story/S3CriticalLineCarrierBundle.lean} ---
      line/off-line carrier contrast: maximal Hopf circle and unimodular
      tangent only on $\Re s=\tfrac12$, shrunk fiber and hyperbolic tangent
      sorting off-line, Path C via rolling on the line;
\item \leanfile{Hqiv/Story/S3HopfJKUnitCircleZeroReadout.lean} ---
      critical line as $j$--$k$ unit circle; $e^{\pi i j k}$ phase readout;
      amplitude balance $\Leftrightarrow$ $\zeta$-zero (conditional);
\item \leanfile{Hqiv/Story/S3HopfJKEulerPrimeCircleCounting.lean} ---
      Euler polar split on the line; prime phases at shell slots; two-prime
      circle pinning; $\zeta$-zero $\Leftrightarrow$ prime twiddle at slot;
\item \leanfile{Hqiv/Story/S3HarmonicShellZeroCounting.lean} ---
      height $t$ as non-injective cover; harmonic-shell slots $(n,k)$ at arc
      angles $2\pi k/n$; sweep refines under $n\mid m$; countable slot encoding;
\item \leanfile{Hqiv/Story/S3AnalyticStripLift.lean} --- cylinder lift
      $\texttt{stripPointLift}$, $45^\circ$ free coordinate and Hopf radius,
      $2\pi$ periodicity (\texttt{sameStripHeight}), $\sigma$-equator
      $=\tfrac12$, whole-strip twiddle;
\item \leanfile{Hqiv/Story/S3AnalyticStripClassification.lean} ---
      conditional zero classification on the full strip (cancellation,
      Path C/E, rolling recovery on the line, centered model);
\item \leanfile{Hqiv/Story/S3LogPhaseEdge.lean} --- polar decomposition
      and decoupling theorem, log edge
      $\texttt{linePhase}(pq)=\texttt{linePhase}(p)\texttt{linePhase}(q)$,
      prime-log incommensurability over $\mathbb{N}$ and $\mathbb{Z}$,
      two-prime phase height pinning, Goldbach circle
      $pq+(q-N)^2=N^2$, phase-speed cap with diagonal equality case,
      RH $\Rightarrow$ zeros addressed by two prime phases.
\end{itemize}

\appendix
\section{Claim status table}\label{app:status}

\begin{table}[h]
\centering
\small
\begin{tabular}{@{}p{0.52\linewidth}cc@{}}
\toprule
Claim & Lean & RH-hard \\
\midrule
$H_n\le K(n)$, $K$ strictly increasing & yes & no \\
$\langle\mathfrak{so}(3),\Delta_4\rangle=\mathfrak{so}(4)$ & yes & no \\
Shell sum $=\zeta(s)$ for $\Re s>1$ & yes & no \\
Prime Euler product $=\zeta(s)$ & yes & no \\
FE strip closed form~\eqref{eq:fe-strip} & yes & no \\
Bernoulli / even-$\pi$ regional values & yes & no \\
Even $s=2k$: $\mathcal{S}(s)=0$, $\mathcal{C}(s)=(-1)^k$ & yes & no \\
Odd $s=2k+1$: $\mathcal{C}(s)=0$ & yes & no \\
$1+\alpha/3=6/5$ harmonic-third projection & yes & no \\
Nontrivial zeros $\notin$ negative integers & yes & no \\
Nontrivial zeros in open strip $0<\Re s<1$ & yes & no \\
$\zeta=h\cdot(2\sigma-1)/\sqrt2$ off $\sigma=\tfrac12$ & yes & no \\
Readout factors through $\Re s$ only & yes & no \\
No global $\zeta=\texttt{exactTwiddleReadout}$ & yes & no \\
FE pair at nontrivial zeros & yes & no \\
Quadruplet zeros (with $\zeta(\overline{s})=\overline{\zeta(s)}$) & yes & no \\
Midpoint tangent slope $N/(p+q)=\tfrac12$ for pairs & yes & no \\
Hopf holonomy support for midpoint pairs & yes & no \\
Bridge payload $\Leftrightarrow$ midpoint Goldbach & yes & no \\
Bridge at $2$ $\Leftrightarrow$ \texttt{GoldbachParity} & yes & no \\
\texttt{SO8ProjectedHalfSlopeBridge}\,$2$ $\Leftrightarrow$ RH $\wedge$ Goldbach & yes & no \\
$\mathfrak{so}(4,2)$ dilation slot count $15=6+4+4+1$ & yes & no \\
Backward weights $>1$, strict growth, unbounded & yes & no \\
Arrow (2nd law) $\Rightarrow$ no backprojection & yes & no \\
$\lambda_{\mathrm{HQIV}}=0\Leftrightarrow\lambda_{\mathrm{HQIV}}\le0$ & yes & no \\
Vaporization bridge $\Leftrightarrow$ RH (given witness) & yes & no \\
$\Phi$ terms $\ge0$ and summable ($u\ge0$) & yes & no \\
Gaussian factor $=$ lattice weight at $(-t,u^2)$ & yes & no \\
$\Phi>0$, continuous, $\Phi\le Ce^{-3u}$, integrable on $(0,\infty)$ & yes & no \\
$\lVert H_t(z)\rVert\le\int\Phi$ for $t\le0$, $z$ real & yes & no \\
Horizon shell weight $=1$ for all $\tau$ & yes & no \\
$c\lVert H_t(z)\rVert\le\lVert\texttt{HQIVDeformedSum}\rVert$ (anchor region) & yes & no \\
Phase alignment $\cos(\Delta\arg)=1$ at anchor & yes & no \\
\texttt{dbnLambda} $=\inf$ hyperbolicity locus (concrete def.) & yes & no \\
$H_t(0)\ne0$ for $t\le0$ (anchor nondegeneracy) & yes & no \\
Tangent orbit law $T(1-s)=T(s)^{-1}$, $T(\overline{s})=\overline{T(s)}$ & yes & no \\
$\lVert\sin\rVert^2-\lVert\cos\rVert^2=-\cos(\pi\sigma)$ at $\pi s/2$ & yes & no \\
Critical line $\Leftrightarrow$ $\lVert\tan(\pi s/2)\rVert=1$ (strip) & yes & no \\
Strip cylinder lift on $S^3$; $t\sim t+2\pi n$ same point & yes & no \\
Line carrier: $f=0$, $r=1$, unimodular $T$, orbit collapse & yes & no \\
Off-line: $r<1$, tangent sorts half-planes, no quadruplet collapse & yes & no \\
Whole-strip cancellation $f(\sigma)+r(\sigma)(\cos t+\sin t)=0$ & yes & yes \\
Path C/E classification under lift $+$ bridge & yes & yes \\
Tangent orbit collapse $T^{-1}=\overline{T}$ $\Leftrightarrow$ line & yes & no \\
Equator factor $=0$ $\Leftrightarrow$ unimodular tangent & yes & no \\
Euler factors nonvanishing; no cancellation for $\Re s>1$ & yes & no \\
Zero $\Rightarrow$ product breakdown ($\Re s<1$, incl.\ boundary) & yes & no \\
Tame spectrum: $\sum\Lambda(n)n^{-s}=-\zeta'/\zeta$ for $\Re s>1$ & yes & no \\
Resonance: $\lVert1/\zeta\rVert\to\infty$ at every zero $\rho\neq1$ & yes & no \\
Spectral-cancellation profile of nontrivial zeros & yes & no \\
Single prime lines/finite truncations $\neq0$ for $\Re s>0$ & yes & no \\
Two-channel dichotomy at every nontrivial zero & yes & no \\
\texttt{InteriorAssemblyNonzero} $\Leftrightarrow$ RH (equivalence) & yes & no \\
FE pairing $\lVert n^{-s}\rVert\lVert n^{-(1-s)}\rVert=1/n$ (all $s$) & yes & no \\
Locators: geometric mean / $\lVert n^{-s}\rVert^2=1/n$ $\Leftrightarrow$ line & yes & no \\
RH $\Leftrightarrow$ square-root spectral weights at zeros & yes & no \\
Mirror fixed set $=$ axis line; FE closure of zero set & yes & no \\
FE-partner collinearity $\Rightarrow$ line is $\Re=\tfrac12$ & yes & no \\
Unique vertical axis (abstract strip sets) & yes & no \\
Unique vertical axis for $\zeta$ zeros (given a zero) & yes & no \\
Weight locus globally $=$ straight line $\Re=\tfrac12$ & yes & no \\
Weights $n^{-2\sigma}$ consistent at every $\sigma$ (collectivity blind) & yes & no \\
Same-height partner $1-\overline{s}$; merge $\Leftrightarrow$ line & yes & no \\
RH $\Leftrightarrow$ height pins $\Re$ (FE only) & yes & no \\
RH $\Leftrightarrow$ one zero per height (given Schwarz) & yes & no \\
Polar separation $=-\sqrt2\cdot\texttt{so4CriticalFactor}$ & yes & no \\
Separation norm $=2\,\lvert\sigma-\tfrac12\rvert$; collapse $\Leftrightarrow$ line & yes & no \\
Single-point polar projection $\Leftrightarrow$ RH & yes & no \\
Pair activation at every zero; joint line factorization & yes & no \\
AM--GM joint-weight cap $1/(pq)\ge1/N^2$ on the line & yes & no \\
Zero contains pair holonomy (slope $+$ Hopf certificate) & yes & no \\
Bridge $\Leftrightarrow$ polar collapse $\wedge$ parity & yes & no \\
Frame tower: nesting, flow rigidity, positivity & yes & no \\
Harmonic radius law $\texttt{frameNormSq}=H_N\Leftrightarrow$ line ($N\ge2$) & yes & no \\
RH $\Leftrightarrow$ harmonic tower radii at every zero & yes & no \\
Radius ladder diverges on the line; $H_N\le K(N)$ bound & yes & no \\
Transformer trace ladder $\to\zeta$ on $\Re s>1$ & yes & no \\
HS norm of $D_N$ $=$ frame radius & yes & no \\
$D_N(s)D_N(1-s)=\mathrm{diag}(1/n)$ (all $s$); trace $=H_N$ & yes & no \\
Adjoint law $D_N(1-s)=D_N(s)^{\dagger}\Leftrightarrow$ line & yes & no \\
RH $\Leftrightarrow$ transformer adjoint law at every zero & yes & no \\
Transformer family abelian and normal (saturation) & yes & no \\
$[e_1,e_2,e_4]=-e_3-e_4$, norm$^2=2$ (frozen tables) & yes & no \\
Associator channel $=2(abc)^{-2\sigma}$; positive at zeros & yes & no \\
Channel $=2/(abc)$ $\Leftrightarrow$ line; FE pair $=4/(abc)^2$ & yes & no \\
RH $\Leftrightarrow$ associator torsion weights at zeros & yes & no \\
Goldbach triple channel floor $1/N^3$ on the line & yes & no \\
Torsion $=(abc)^{-\sigma}(-e_3-e_4)$: single S$^7$ point & yes & no \\
FE torsion defect odd; cancels $\Leftrightarrow$ line & yes & no \\
Chirality flip of the defect across the line & yes & no \\
SO(8) orbit coincidence $\Leftrightarrow$ vector collapse $\Leftrightarrow$ line & yes & no \\
RH $\Leftrightarrow$ S$^7$ torsion cancellation at every zero & yes & no \\
Polar decoupling: modulus $\sigma$-only, phase $t$-only & yes & no \\
Prime-log incommensurability ($\mathbb{N}$, $\mathbb{Z}$) & yes & no \\
Two prime phases pin the height exactly & yes & no \\
Goldbach circle $pq+(q-N)^2=N^2$ & yes & no \\
Phase-speed cap $\log(pq)\le2\log N$; diagonal equality & yes & no \\
RH $\Rightarrow$ zeros addressed by two prime phases & yes & no \\
Rodgers--Tao $0\le\texttt{dbnLambda}$ (stated; proof) & no & --- \\
Newman/dBN RH $\Leftrightarrow\texttt{dbnLambda}\le0$ (stated; proof) & no & --- \\
$\texttt{dbnLambda}\le\lambda_{\mathrm{HQIV}}$ (inheritance bound) & no & yes \\
\texttt{DeltaHarmonicCompleteness} $\Rightarrow$ Goldbach & no & yes \\
\texttt{interiorStripH}$\neq0$ at nontrivial zeros off line & no & yes \\
Witness for \texttt{SO8ProjectedHalfSlopeBridge}\,$2$ & no & yes \\
Height $t$ is cover, not injective counter on $S^1$ & yes & no \\
Harmonic-shell slots $(n,k)$ at $2\pi k/n$, encodable & yes & no \\
Arc sweep refines when $n\mid m$ & yes & no \\
$j$--$k$ unit circle $(\cos t,\sin t)$ on Hopf fiber & yes & no \\
$e^{\pi i j k}$ phase readout on circle & yes & no \\
$\zeta=0$ on line $\Leftrightarrow$ $j/\sqrt2+k/\sqrt2=0$ & cond.\ bridge & yes \\
Harmonic-shell slot inherits $j$--$k$ twiddle readout & yes & no \\
$\zeta=0$ on line $\Leftrightarrow$ rolled cover balance & cond.\ bridge & yes \\
Unconditional \texttt{RiemannHypothesis} & no & yes \\
Unconditional \texttt{GoldbachParity} & no & yes \\
\bottomrule
\end{tabular}
\end{table}

\section{Notation}\label{app:notation}

\begin{tabular}{@{}ll@{}}
$\alpha$ & HQIV imprint $3/5$ (\leanid{alpha\_eq\_3\_5}) \\
$H_n$ & harmonic partial sum (\texttt{harmonicPartialSum}) \\
$(n,k)$ & harmonic-shell slot, angle $2\pi k/n$ (\texttt{HarmonicShellSlot}) \\
$K(n)$ & curvature integral (\texttt{curvature\_integral}) \\
$\texttt{so4CriticalFactor}(s)$ & $(2\sigma-1)/\sqrt2$ \\
$\texttt{oddStripChannel}$ & FE assembly~\eqref{eq:fe-strip} \\
$\texttt{interiorStripH}$ & $(\texttt{even}+\texttt{odd})/\texttt{so4CriticalFactor}$ \\
$\mathcal{S},\mathcal{C}$ & $\sin(\pi s/2)$, $\cos(\pi s/2)$ \\
\end{tabular}

\bibliographystyle{amsplain}
\bibliography{../references}

\end{document}
